% =====================================================================
%  Capitolo 14 — Illuminare con uno scopo: il Physically Based
%  Rendering, la mobilità delle luci, il sole, il cielo e la
%  verosimiglianza in Unreal Engine (Game Design, primo anno)
%  Fonte: trascrizione della lezione GD14, terza lezione del corso di
%  Level Design del primo anno (docente: Stefano Zocchi, lo stesso dei
%  capitoli 8 e 11), di cui è la diretta prosecuzione: i task rivisti
%  in apertura sono le stanze inventate assegnate nel capitolo 11.
%  Contenuti, nell'ordine della lezione: la chiacchierata d'apertura su
%  GTA 6 e su Fortnite (il token, il gioco d'azzardo, la piattaforma);
%  i due pacchetti della lezione (Medieval Village e Scan Labs
%  Interiors); la revisione delle stanze inventate (lo studente dei
%  booleani, Valerio, Alberto, Francesco, Federica, Antonino,
%  Francesca) con le digressioni sulla coerenza strutturale, sulle
%  reference dal vero e sul manichino per le proporzioni; il ripasso
%  delle viste ortografiche, delle modalità di visualizzazione, del
%  LOD del terreno, delle gerarchie e della differenza tra figli e
%  componenti; il Physically Based Rendering (texture e materiale,
%  specular highlight, albedo, roughness, specular, la mappa della
%  roughness); i quattro tipi di luce e la mobilità (Static,
%  Stationary, Movable), il baking e le lightmap, Lumen, Build
%  Lighting Only; la directional light e il SunSky di Medieval
%  Village; la notte come giorno tinto di blu; le Cascade Shadow Maps;
%  la Sky Light e il colore delle ombre; la point light e i tre esempi
%  di Lies of P (il barile, la lampada senza ombra, la luce motivata
%  del corridoio); la pausa (ancora GTA, la monetizzazione, il finto
%  MMO; la domanda sul posizionamento degli oggetti via codice, le
%  spline e i geometry node); le luci che Unreal spegne da solo e il
%  Cast Shadows; la spot light; Medieval Village giocabile; la mappa
%  di Scan Labs letta insieme (fuori campo non si costruisce, il
%  chiaroscuro, le reflection capture, il pizzico di caos, le tre
%  regole dell'osservazione, la luce rimbalzata, gli artefatti di
%  Lumen, Lumen e il ray tracing); i materiali proiettati dalle luci;
%  il task (una stanza con luci motivate); le domande sul task; la
%  lezione spostata e la chiusura.
%  Integrazioni dal libro "An Architectural Approach to Level Design"
%  (Christopher W. Totten), segnalate come approfondimenti.
%
%  Pacchetti richiesti nel preambolo del libro (già presenti nel
%  main.tex): \usepackage{booktabs}, \usepackage{amsmath},
%  \usepackage{enumitem}, \usepackage{tikz} con
%    \usetikzlibrary{shapes.geometric,shapes.symbols,shapes.misc,
%      arrows.meta,positioning,calc,fit}
%  (le librerie vanno caricate nel preambolo, non nel capitolo).
%  Tutte le etichette di questo capitolo hanno il prefisso "gd14:"
%  (gd14:sec:..., gd14:fig:..., gd14:tab:...) e tutti gli stili TikZ
%  il prefisso "gd14", per non collidere con gli altri capitoli.
%
%  Nota sulla fonte: la trascrizione automatica è buona nelle parti
%  spiegate alla classe (la revisione, il PBR, le luci, la lettura
%  della mappa, il task), ma molto rovinata nelle due pause, nei
%  momenti di lavoro individuale e quando il docente parla a bassa
%  voce con un singolo studente. Quei tratti sono stati trascritti come
%  frasi senza senso o ripetute in loop ("ti piacciono", "sì, sì",
%  "mi fa", "non rispondi?", "ho fatto un po' di lavoro", "non lo so",
%  "se non prendo niente", "mi vuoi dire?", "una tecnica che ha messo a
%  punto"), con lunghi brani captati da una fonte esterna (frasi sulla
%  "signora Gibbardi", sulla "signora Mazzoni", sulla "città-stato",
%  sulla "dottrina sul mercato del lavoro", su "Isabella Santander",
%  sui "funzionari della società", "la migliore dal 2005", "il numero
%  di persone è molto più basso dopo il raggio di persone", "il raggio
%  di vetro", "Vittorio Prodi"): sono artefatti, omessi. Si omettono
%  anche gli scambi privati non legati alla lezione (un torneo di un
%  picchiaduro e un amico del docente, il caricamento del progetto per
%  una studentessa, la riparazione dei materiali sul computer di uno
%  studente, il recupero di una consegna di uno studente lavoratore).
%  Molti nomi sono storpiati ("Arial", "Arreal", "a Real", "un Real",
%  "Harry Potter" per Unreal; "l'humen", "l'UMEN", "lume" per Lumen;
%  "PDR" per PBR; "Raphness" per roughness; "albido" per albedo;
%  "Skylite" per Sky Light; "directionalite", "detection light",
%  "directional line" per directional light; "pointera" per point
%  light; "Scarlabs", "Scala", "Scanlabs" per Scan Labs; "Medival
%  Village", "Medieval Town", "Medieval City" per Medieval Village;
%  "IsoFP", "Lise of P", "Eyes of Pie", "Light of P", "Light of
%  Python", "L'ISO Speech" per Lies of P; "il Diver 2" per Helldivers
%  2; "Aina in barbecue" per ENA: Dream BBQ; "GTH 6" per GTA 6;
%  "refreshing" per ray tracing; "Silent Till Have" per Silent Hill):
%  dove il riferimento è ricostruibile con certezza lo si riporta nella
%  forma corretta; dove non lo è, lo si segnala in nota invece di
%  indovinare. Le figure e le tabelle sono ricostruzioni schematiche di
%  ciò che il docente mostra e descrive, non le slide originali né
%  screenshot.
% =====================================================================

\chapter[Luci, PBR e verosimiglianza]{Illuminare con uno scopo: il PBR,
la mobilità delle luci e la verosimiglianza}
\label{gd14:cap:luci}

Il capitolo~\ref{gd11:cap:blockout-luci} aveva dato agli studenti il
\textbf{vocabolario} del level designer: blockout, coordinate, gerarchie, pivot,
mesh, e una prima carrellata delle luci di Unreal, chiusa con il task di
inventare \textbf{una stanza}. Questa lezione, tenuta dallo stesso docente
(Stefano Zocchi), è la terza del corso di \textbf{Level Design} del primo anno,
ed è dedicata quasi interamente alle \textbf{luci}: ``oggi parliamo di luci,
oggi approfondiamo un attimo il discorso delle luci''. Non è l'unica: il docente
annuncia che, dato che ``c'è un sacco di roba di cui parlare'', le lezioni sulle
luci saranno \textbf{quattro}.

La lezione ha tre anime. C'è una parte \textbf{tecnica}, per capire ``cosa sta
succedendo dietro le quinte dell'engine'': il \emph{Physically Based Rendering},
che spiega perché un materiale reagisce alla luce in un certo modo, e la
\textbf{mobilità} delle luci (statiche, stazionarie, mobili), che spiega perché
una modifica a volte ``non fa niente''. C'è una parte \textbf{artistica}, che è
il vero cuore della lezione: la differenza tra \textbf{realismo} e
\textbf{verosimiglianza}, mostrata con tre esempi tratti da \emph{Lies of P}, e
l'idea che le luci di un videogioco, come quelle di un set cinematografico, si
piazzano \textbf{per uno scopo}, non ``dentro la lampada''. E c'è una parte di
\textbf{osservazione}: due progetti professionali (\emph{Medieval Village} e una
mappa di \emph{Scan Labs}) da aprire, esplorare e smontare, e le \textbf{tre
regole dell'osservazione} che il docente consegna alla classe.

Come quasi sempre, la lezione comincia prima della lezione, con una lunga
chiacchierata sull'uscita di \emph{GTA 6} e sulle ultime trovate di
\emph{Fortnite}: il capitolo la riporta, perché dice molto su come un
professionista guarda l'industria.

La struttura della lezione, e quindi del capitolo, è questa:
\begin{enumerate}
  \item la chiacchierata d'apertura su \emph{GTA 6} e su \emph{Fortnite}
        (sezione~\ref{gd14:sec:gta});
  \item i due pacchetti della lezione, \emph{Medieval Village} e \emph{Scan Labs
        Interiors}, e la differenza tra un progetto e un pacchetto di asset
        (sezione~\ref{gd14:sec:pacchetti});
  \item la \textbf{revisione delle stanze inventate}, con le digressioni sulla
        coerenza strutturale, sulle reference dal vero e sul manichino per le
        proporzioni (sezione~\ref{gd14:sec:revisione});
  \item il \textbf{ripasso}: viste ortografiche, modalità di visualizzazione,
        il LOD del terreno, l'ottimizzazione, le gerarchie e la differenza tra
        figli e componenti (sezione~\ref{gd14:sec:ripasso});
  \item il \textbf{Physically Based Rendering}: texture e materiale, albedo,
        roughness, specular (sezione~\ref{gd14:sec:pbr});
  \item le luci e la loro \textbf{mobilità}: Static, Stationary, Movable, il
        baking, le lightmap e Lumen (sezione~\ref{gd14:sec:mobilita});
  \item la \textbf{directional light}, il SunSky e la notte
        (sezione~\ref{gd14:sec:directional}); la \textbf{Sky Light} e il colore
        delle ombre (sezione~\ref{gd14:sec:skylight});
  \item la \textbf{point light} e la \textbf{verosimiglianza}: i tre esempi di
        \emph{Lies of P} (sezione~\ref{gd14:sec:liesofp});
  \item la pausa: ancora \emph{GTA}, e la domanda sul posizionamento degli
        oggetti via codice (sezione~\ref{gd14:sec:pausa});
  \item le luci che Unreal \textbf{spegne da solo} e il \textbf{Cast Shadows}
        (sezione~\ref{gd14:sec:troppe-luci}); la \textbf{spot light}
        (sezione~\ref{gd14:sec:spot});
  \item \emph{Medieval Village} giocabile e la \textbf{lettura} della mappa di
        Scan Labs: fuori campo non si costruisce, il chiaroscuro, il pizzico di
        caos, le tre regole dell'osservazione, la luce rimbalzata e gli
        artefatti di Lumen (sezioni~\ref{gd14:sec:play}--\ref{gd14:sec:lumen});
  \item i materiali proiettati dalle luci (sezione~\ref{gd14:sec:lightfn});
  \item il \textbf{task}, le domande e la chiusura
        (sezioni~\ref{gd14:sec:task}--\ref{gd14:sec:chiusura}).
\end{enumerate}

% ---------------------------------------------------------------------
\section{Prima della lezione: GTA 6 e Fortnite}
\label{gd14:sec:gta}

La registrazione parte a discussione già avviata: è la settimana dell'uscita di
\emph{GTA 6}, e uno studente ha chiesto al docente che cosa ne pensi. Non è una
digressione da poco: è un piccolo esempio di come un professionista legge un
evento dell'industria, ed è per questo che il capitolo la riporta.

\subsection{Un successo, ma abbastanza?}

\paragraph{La domanda.} Rockstar ha speso ``una quantità allucinante di soldi''
su \emph{GTA 6}, soprattutto per lo sviluppo, ma anche per il
\textbf{marketing}: il solo marketing di \emph{GTA 6}, dice il docente, ``costa
dieci volte lo sviluppo di un videogioco medio''. ``Hanno praticamente colonizzato
Miami'', hanno organizzato una serie di eventi speciali: l'uscita è stata
trattata ``non come il lancio di un videogioco'', ma come il lancio del prodotto
di intrattenimento più costoso della storia. Che sarà il videogioco più venduto
della storia, ``su quello non ci piove''; che sarà un successo enorme di critica
e di pubblico è ``inevitabile''. La vera domanda è un'altra: \textbf{basterà}? Basterà
a ripagare quello che ha speso Rockstar e, soprattutto, a creare quello che
Rockstar si aspetta, per il gioco ``considerato quello che salverà l'industria
dei videogiochi''?

\paragraph{Una risposta che arriverà tra anni.} ``Hai posto una domanda la cui
risposta non avremo fra un mese, non avremo fra un anno.'' Il motivo è che lo
scopo di Rockstar è fare di \emph{GTA 6} \textbf{la piattaforma di lancio del
nuovo GTA Online}. Se sarà stata una buona mossa, o una ``mossa alla
\emph{Destiny}'', dove alla fine si è dovuto concludere ``abbiamo perso un sacco
di soldi'', lo si saprà tirando le somme tra tre o quattro anni. ``Adesso sono
tutti in cerca del facile: è un successo gigantesco, è un fallimento clamoroso.''
Ognuno tirerà acqua al suo mulino; Rockstar dirà di sicuro che è un successo
clamoroso, ``perché sarà un successo clamoroso'', ma non si saprà se è il
successo \emph{che si aspettavano} finché non si sommeranno i ricavi di
\emph{GTA 6} e di \emph{GTA Online}. Anche se \emph{GTA 6} andasse in perdita, se
\emph{GTA Online} diventasse ``una roba che gli fa fare cinquanta milioni di
dollari al mese'', il discorso sarebbe quanto in là riesce a tirare, per rientrare
dei costi e continuare lo sviluppo.

\paragraph{Il punto di vista del professionista.} Il docente non ha una
PlayStation~5, ``quindi non ci posso giocare a prescindere''. Lo incuriosisce
però ``lato professionale'' vedere \textbf{che cosa si può ottenere in un
videogioco in cui hai una quantità di soldi infinita} per fare quello che vuoi. E
a uno studente che teme che dieci anni di sviluppo non siano un buon segnale
risponde di no: ``hanno così tanti soldi, così tanto personale e
\textbf{l'esperienza per lavorare} con così tanti soldi e così tanto personale,
che sarà un gran bel gioco. Terra-terra, sarà un gran bel videogioco. Vedremo se
è il videogioco di cui avevano bisogno.''

\paragraph{Solo Rockstar può fare quello che fa Rockstar.} Lato industria, ci sarà
``un bel po' di gente che prova a imitare \emph{GTA 6}'', come c'è stata gente che
provava a imitare \emph{GTA 5} e \emph{GTA 4}, e fallirà clamorosamente. Il motivo
è pratico: altri studi potrebbero avvicinarsi a un budget comparabile, ma
\textbf{nessun altro studio può garantire il ritorno} che ha \emph{GTA}, ``anche
solo per il nome''. \emph{GTA} è una delle rarissime serie a cui gioca anche
\textbf{chi non gioca ai videogiochi}.

\paragraph{La calma piatta.} Un consiglio pratico: ``godetevi la calma piatta'',
perché quando esce \emph{GTA} nessun altro osa far uscire un gioco per un mese;
``e state lontani dai social'', perché per tre settimane parleranno solo di
quello.

\subsection{Fortnite: il token e il gioco d'azzardo}

La conversazione passa a \emph{Fortnite}. Qual è la grande strategia di Epic per
rilanciare un gioco che ``sta iniziando a calare''? Uno studente racconta del
\textbf{token} che i giocatori danno la caccia a trovare, e di conoscenti che
``perdono ore di sonno'' per trovarlo. La risposta del docente è secca:
``\textbf{il gioco d'azzardo}. Letteralmente, il gioco d'azzardo.'' A un certo
punto, ``chiediti se ti stai divertendo o se sei vittima di ludopatia''.

\paragraph{Un simulatore di gioco d'azzardo.} Per il docente \emph{Fortnite} è già
di base ``un simulatore del gioco d'azzardo'', a livello di design ma anche con
lo \textbf{store}, con ``tutte le robe in rotazione'' e i suoi ``trucchi
squalidi'' (ricorda che Epic è stata anche multata dall'Unione Europea per i
suoi ``trucchetti psicologici''\footnote{Così nella lezione: è un'affermazione
del docente, riportata come tale.}). L'analogia per far capire il token: ``è come
se Valve dicesse, in \emph{Deadlock}: in ogni partita hai una possibilità su
diecimila che, se vinci, ti diamo cinquanta euro di buono per Steam''. ``Ok, va
bene: è letteralmente gioco d'azzardo. Io non le sopporto, queste robe.''

\paragraph{Perché è così popolare.} Il docente cerca di giocare ai free-to-play
quando iniziano a diventare popolari, ``per capire un attimo cosa sta
succedendo''. La sua conclusione su \emph{Fortnite}: ``un gioco decente, ma niente
di speciale''. Ciò che lo rende così popolare è che \textbf{ogni mese, o
addirittura ogni due settimane, aggiungono qualcosa}: una rotazione continua di
nuove meccaniche, nuovi personaggi, nuove skin, eventi speciali, collaborazioni,
mappe a tema. Questa rotazione ``sovrasta qualunque mancanza del gioco in sé come
gioco'', al prezzo di ``spremere come delle spugne gli sviluppatori'', che devono
produrre una quantità di contenuti ``fuori da ogni comprensione''.

\subsection{Non un gioco: una piattaforma}

\paragraph{I Funko Pop.} ``Il loro interesse non è mai stato fare un vero gioco,
il loro interesse è stato attirare la gente.'' La formula del docente:
\begin{quote}
\emph{Fortnite} sta ai videogiochi come i Funko Pop stanno alla scultura.
\end{quote}
Lo scopo è sempre stato creare uno \textbf{spazio sociale}: la gente non gioca a
\emph{Fortnite} per vincere, ma per passare il tempo con gli amici (o, nell'ultima
incarnazione, per cercare il token). È un \textbf{terzo spazio}, ``strapieno di
roba che ti vendono'' e di collaborazioni con altre aziende (\emph{The Amazing
Digital Circus}, \emph{Guerre Stellari}, ``chi più ne ha più ne metta''): un
generico ``spazio di divertimento corporate con tutti i tuoi personaggi
preferiti'', dove puoi vedere Goku che balla sul cadavere di un altro personaggio,
``quel tipo di cosa che una volta si faceva con i disegni in Paint messi su
internet''. Il gioco ha comunque ``un suo perché'' nella casualità, nel calderone
di meccaniche e veicoli ``messi un po' sparsi a caso, che tanto la settimana
prossima cambiano''.

\paragraph{Il modello Roblox.} ``\emph{Fortnite} non è un gioco, è una
piattaforma. O meglio: vogliono che sia una piattaforma.'' Nominalmente è un
gioco, ma in pratica vuole essere \textbf{Roblox}. E che cos'è Roblox? ``L'Amazon
dei videogiochi'', con il suo engine proprio. ``Nessuno vuole vendere al mercato:
tutti vogliono \textbf{essere il mercato}, dove prendi una percentuale su tutto
quello che viene venduto al tuo interno.'' Per questo Epic ha spinto molto sulle
\textbf{isole}, cioè sui contenuti prodotti dal pubblico invece che dagli
sviluppatori. ``Non sta andando'', chiude il docente, e si passa alla lezione.

% ---------------------------------------------------------------------
\section{I pacchetti di oggi: Medieval Village e Scan Labs}
\label{gd14:sec:pacchetti}

\paragraph{Scaricare subito.} Per la lezione servono dei file, e il docente chiede
di scaricarli subito ``così non perdiamo tempo''. Il wi-fi dell'aula non regge
(``saluti al wi-fi in classe; io sono connesso via cavo, quindi non è troppo un mio
problema''), così il docente fa girare una \textbf{chiavetta USB} con i file:
``ci sono due file, prendeteveli entrambi, per favore; ignorate gli altri''. Se
il primo si è già scaricato dal link, dalla chiavetta basta prendere il secondo.
Il secondo pacchetto è sulla chiavetta solo perché il docente lo ha preparato la
sera prima, ed era più veloce passarlo a mano in classe; l'altro era già online.
A una studentessa collegata da casa (Federica) il docente lo carica durante la
pausa, così può scaricarlo.

\paragraph{Due cose diverse.} I due pacchetti \textbf{non vanno trattati allo stesso
modo}, ed è l'errore che uno studente ha già fatto:
\begin{itemize}
  \item \textbf{Scan Labs Interiors}\footnote{Nella trascrizione ``Scarlabs'',
        ``Scan Labs'', ``Scala'', ``Scanlabs''; la mappa al suo interno si chiama
        \emph{Interiors Volume 1}.} è un \textbf{pacchetto di asset} già
        importati, come quelli visti la volta precedente: va messo
        \textbf{dentro la cartella \texttt{Content}} del proprio progetto, come
        si è fatto per lo Starter Content
        (sezione~\ref{gd11:sec:starter}). Chi usa il launcher di Unreal può
        cliccare sul progetto e scegliere \emph{Mostra nella cartella}: si apre
        la cartella del progetto, ci si butta dentro il pacchetto ``e siete a
        posto''. Ricordarsi di \textbf{decomprimerlo}. Non serve un progetto
        nuovo: va bene anche uno già usato, creando un livello nuovo.
  \item \textbf{Medieval Village} è un \textbf{progetto} intero: ``non va messo
        dentro un altro progetto''. È proprio un file di Unreal: ci si fa doppio
        clic e lo si apre. Unreal chiederà con che versione aprirlo, perché è
        stato fatto con una versione precedente (``in 5.6, mi sembra'', o in
        5.7): si sceglie la \textbf{5.8}, ``funziona comunque benissimo''. Allo
        studente che lo aveva copiato dentro il proprio progetto: ``cancella
        quello che hai messo dentro''.
\end{itemize}

\paragraph{Perché Medieval Village.} \emph{Medieval Village} è il lavoro di un
\textbf{level designer italiano}\footnote{Più avanti il docente ne pronuncia il
nome mentre scorre la schermata iniziale; nella trascrizione è reso come ``Marco
Maria Rossi'', che non è verificabile.}: un villaggio medievale esplorabile,
``a un livello di qualità da gioco AAA''. Il docente è chiaro: ``è ben oltre
quello che siete in grado di fare, è veramente fatto da una persona
espertissima''. Ma proprio per questo è utile: è un \textbf{punto di riferimento}
di come si fa questo lavoro a livello professionale, e un progetto in cui si può
\textbf{mettere le mani} e guardare dietro le quinte. La raccomandazione, comunque,
è di aspettarsi di non capirci niente: si tornerà sul punto
(sezione~\ref{gd14:sec:directional}).

% ---------------------------------------------------------------------
\section{La revisione delle stanze inventate}
\label{gd14:sec:revisione}

Mentre la chiavetta gira e Unreal si apre (``questo è un computer nuovo, forse ci
mette poco''), si guardano i task: le \textbf{stanze inventate} assegnate nella
lezione precedente (sezione~\ref{gd11:sec:task}). Uno studente aveva scritto al
docente di avere avuto problemi a importare i materiali, e non li ha risolti: lo
si vedrà dopo, a parte. La cartella non si lascia ordinare per nome (``fantastico''),
quindi il docente va ``in ordine al primo che vede''.

\subsection{Lo studente dei booleani: spazi, non forme a caso}

Il primo lavoro è fatto di forme ottenute con operazioni \textbf{booleane}, cioè
lavorando con la \textbf{modellazione} dentro Unreal. Lo studente ha già seguito
corsi extra di Unreal, e si vede: ``forme a caso realizzate con una bella
coscienza tecnica, per carità'', ma la \emph{Modeling Mode} si vedrà più avanti,
``quindi non ce n'era bisogno''.

\paragraph{Creare un piano e seguirlo.} Il richiamo vale per tutti. Le prime
lezioni servono a mettere le mani sull'engine, ma il docente vuole anche
sviluppare ``la capacità di creare un piano e seguirlo per creare qualcosa di
\textbf{strutturale}''. ``Io vi ho chiesto di creare una stanza, e queste sono
forme a caso.'' L'obiettivo personale che il docente dà allo studente:
\textbf{sviluppare coerenza strutturale}, creare \textbf{spazi}, ``non cose messe a
caso per sapere che lo sai fare''. Uno degli obiettivi del corso è sviluppare un
\textbf{senso estetico}, un senso di struttura: il primo anno è tutto dedicato
all'aspetto estetico (luci, materiali, ambienti).

\paragraph{Niente avanguardia, per ora.} Le prime lezioni servono a prendere la
mano, e il docente non è ``esageratamente rigido''. Ma da ora si comincia a
indirizzare l'impegno verso risultati ``che ipoteticamente potremmo vedere in un
videogioco''. E senza la scusa che qualcuno tira fuori ogni tanto: ``ah, ma hai
giocato a \emph{ENA: Dream BBQ}?\footnote{Nella trascrizione ``Aina in
barbecue''.}'', o ad altri giochi super indie con visuali strane come
\emph{Cruelty Squad}. ``Non prendiamoci in giro: cerchiamo di fare ambienti logici
e \textbf{lasciamo l'avanguardia a quando abbiamo finito le basi}.'' E ancora: ``non
voglio più vedere cose astratte dadaistiche da ora in poi''; lo scopo è
concentrarsi su ``qualcosa di definito e definibile'', ``qualcosa che potrebbe
essere una scena di un videogioco''.

\paragraph{Usare una reference.} Il consiglio pratico: ``usa reference, se ne hai
bisogno. Guarda una fotografia di una stanza e cerca di ricostruire quella
stanza.'' È anche uno dei motivi per cui il docente fa girare \emph{Medieval
Village}: avere sotto gli occhi un lavoro professionale.

\subsection{Valerio: la casa sulla sedia}

Valerio aveva scritto al docente di avere avuto grossi problemi con l'engine: i
materiali non si importavano, e ha dovuto usare dei \textbf{colori} al posto delle
texture. ``Tuttavia direi che hai ottenuto un ottimo risultato, a prescindere.''
L'idea è una \textbf{casa costruita su una sedia}: era un esercizio di fantasia,
``quindi ci sta''; molto moderna, con gli spazi interni ben organizzati. C'è
anche la pianta del pacchetto, i cui materiali si sono rotti con l'importazione
sbagliata: Valerio ``l'ha vista così e l'ha trasformata in un esempio di arte
moderna. Mi piace come pensi.''

\subsection{Alberto: verso il realistico}

Alberto si sposta ``un pochino più sul realistico, o sul tendente realistico'',
``e ci sta anche qua''. Ha costruito a mano anche alcuni elementi, con
impegno. ``Casa, perfetto, niente da dire. Hai fatto un gran bel lavoro.''

\subsection{Francesco: l'aula, e il pavimento che tutti riconoscono}

Francesco (Muratore) ha \textbf{rifatto l'aula} in cui si tiene la lezione. ``Molto
carina come idea.'' Il docente apprezza l'attenzione ai dettagli delle
\textbf{prese di corrente sui banchi} e il lavoro sul materiale del pavimento, un
legno dello Starter Content: ``a furia di insegnare e di vedere gente che usa i
pacchetti di Unreal di base, ormai lo conosco a memoria quel pavimento. Appena lo
vedo in un gioco, mi dico: \emph{Starter Content}.''

\paragraph{Mai uno a uno.} Quando si rifanno ambienti reali in un videogioco, ``non
si fanno quasi mai uno a uno'': c'è sempre un certo livello di
\textbf{reinterpretazione degli spazi}. Perfino giochi come \emph{Assassin's
Creed}, che si vantano di aver replicato città nel modo più realistico possibile,
``un minimo di libertà se la prendono'': a volte per necessità tecniche, a volte
perché di una certa parete non ci sono foto e bisogna improvvisare.

\subsection{Federica: la vecchia stanza, e il manichino}

Federica Ferrara segue da casa (``mi senti?'' ``Forte e chiaro.''). Ha fatto un
interno di casa con una \textbf{statuina con due materiali diversi}, una
libreria con i libri, la scrivania, la stampante, perfino la PlayStation~5. ``Hai
un buon occhio per i dettagli, hai un buon occhio per la ricostruzione degli
ambienti.'' È la sua \textbf{vecchia stanza}. ``Super attenzione ai dettagli,
onestamente ben oltre quello che mi aspetto dalle vostre task, ma mi fa molto
piacere vederlo.''

\paragraph{Il manichino per la scala.} In mezzo alla stanza c'è il
\textbf{manichino di Unreal}: Federica lo ha messo per dare la proporzione di
tutto. Il docente approva: è ``un classico trucco della modellazione 3D per
mantenere le dimensioni corrette''. Si prende un modello 3D del proprio
personaggio, lo si mette nella scena e \textbf{ci si costruisce tutto attorno},
usandolo come punto di riferimento.

\subsection{La lezione delle reference}

Il docente fa notare una cosa: sia Francesco sia Federica hanno usato come
riferimento \textbf{stanze che conoscono} (l'aula, la propria vecchia camera). E
``nel momento in cui si usa una reference dalla realtà, riusciamo a inserire molti
più dettagli''.
\begin{quote}
Quando si lavora a livello artistico, è un buon modo \textbf{partire dalla realtà e
modificarla}, piuttosto che partire direttamente dalla fantasia.
\end{quote}
Chi si mette a imparare a disegnare se lo sente dire per primo: \textbf{usate le
reference}, partite da fotografie. ``Non ricalcate le fotografie, ma tenetele
sott'occhio per dare un'occhiata a tutti i dettagli.'' La cosa tornerà anche nel
mondo dei videogiochi.

\paragraph{Niente asset da Fab.} Una piccola nota: una studentessa aveva chiesto su
Discord il permesso di aggiungere prop scaricati da \textbf{Fab} (e in uno degli
screenshot c'era un gatto). Il docente ha detto di no, ``perché voglio che partiate
tutti dallo \textbf{stesso punto di riferimento}''. La regola resta valida anche per
il task successivo (sezione~\ref{gd14:sec:domande}).

\paragraph{Il tasto G.} E un promemoria: premendo \textbf{G} nell'editor si tolgono
tutte le icone dalla visualizzazione, così gli screenshot sono puliti.

\subsection{Antonino: Helldivers 2, e il tappeto di muschio}

Antonino ``si sposta molto più sul fantascientifico'': voleva ispirarsi a
\emph{Helldivers 2}.\footnote{Nella trascrizione ``il Diver 2''.} Un problema
tecnico: un materiale \textbf{trasparente} non riceve l'illuminazione come gli
altri. Due cose che funzionano:
\begin{itemize}
  \item il \textbf{tappeto} che riempie lo spazio, fatto con un materiale d'erba (è
        muschio): ``se non stai a farci il pelo nell'uovo, è leggibile come
        tappeto''. Ed è il criterio generale: ``se qualcosa è
        \textbf{leggibile} come quello che deve essere inteso, allora stiamo
        andando nella giusta direzione'';
  \item ``un piccolo esempio di \textbf{caos}'': le sedie spostate ad angoli
        casuali, i tablet spostati ad angoli casuali. Queste cose rendono gli
        ambienti ``un po' più vissuti, leggermente più caotici, e quel pizzico
        di caos li rende meno artificiali, meno artificiosi'' (lo stesso
        principio emerso nel capitolo~\ref{gd11:cap:blockout-luci}).
\end{itemize}

\subsection{Francesca: lo chalet fatto di tubi}

Francesca, ``ultima ma non ultima'', ha fatto uno \textbf{chalet di montagna}
fatto\ldots\ con dei tubi. ``L'intenzione c'è ed è leggibile, quindi ottimo.'' Il
perché è semplice: non c'erano i tronchi.

\paragraph{Arrangiarsi con quello che si ha.} ``Questa è una situazione in cui vi
troverete di sicuro: volete fare qualcosa ma non avete i mezzi per farlo.'' In
questo caso va benissimo, perché il tipo di ambiente si legge. Ma ci si può
chiedere: i tronchi non li posso fare, perché non ho il materiale necessario;
allora facciamo delle \textbf{assi di legno}, usando il materiale del pavimento.
Altra osservazione: pavimento e pareti hanno lo \textbf{stesso aspetto}, ``e questo
crea una roba molto piatta''.

\subsection{Lo scopo degli esercizi}

``Come dicevo in precedenza, lo scopo di questi esercizi è farvi prendere la mano
con Unreal e farvi sviluppare l'istinto del \emph{mi viene in mente, lo faccio}'',
senza fermarsi a cercare dove sono i controlli. Uno studente non ha fatto
l'esercizio (la settimana precedente è stato male): ``cerca di recuperarlo, perché
comunque voglio vederli tutti, questi esercizi''.

% ---------------------------------------------------------------------
\section{Ripasso: viste, modalità di visualizzazione e gerarchie}
\label{gd14:sec:ripasso}

Si torna alla lezione riprendendo i concetti spiegati la volta precedente, che il
docente aveva messo anche all'inizio delle slide di oggi. Prima un trucco già
detto: tenendo premuto il \textbf{tasto destro} e girando la \textbf{rotella} del
mouse si aumenta o diminuisce la \textbf{velocità della telecamera}. Tutte le
scorciatoie sono nella \textbf{prima slide} con le scorciatoie di Unreal.

\subsection{Perspective, Top, Bottom, Left, Right}

C'è una cosa che il docente non aveva spiegato. In alto a sinistra nella viewport,
dove c'è scritto \textbf{Perspective}, si può cambiare la \textbf{proiezione} della
vista (tabella~\ref{gd14:tab:viste}).
\begin{itemize}
  \item \textbf{Perspective} è la vista di base: la normale telecamera ``del mondo
        che ci circonda, che simula il modo in cui vedono i nostri occhi''.
  \item \textbf{Top} mostra il mondo di gioco \textbf{dall'alto}, come una mappa.
        Passando a questa vista, Unreal \textbf{spegne le luci}: la modalità
        diventa automaticamente \emph{Unlit}, ``quindi è molto facilmente
        leggibile come una mappa''. Lo stesso vale per \textbf{Bottom} (dal basso
        verso l'alto), \textbf{Left}, \textbf{Right} e le altre.
\end{itemize}
Sono viste \textbf{completamente bidimensionali} (ortografiche), utili per valutare
i vari lati di quello che si sta facendo. Se si sta piazzando qualcosa dall'alto,
per esempio la disposizione di una mappa, in \emph{Top} si spostano gli oggetti
comodamente. Tornando a \emph{Perspective}, la telecamera torna dov'era, con tutti
i cambiamenti fatti.

\begin{table}[htbp]
\centering
\small
\begin{tabular}{@{}p{3.0cm}p{4.6cm}p{5.0cm}@{}}
\toprule
\textbf{Vista} & \textbf{Che cosa mostra} & \textbf{Quando serve} \\
\midrule
Perspective & la telecamera ``come i nostri occhi'', in prospettiva &
  il lavoro di tutti i giorni, gli screenshot \\
\addlinespace
Top / Bottom & il livello dall'alto o dal basso, senza prospettiva, in
  \emph{Unlit} & leggere il livello come una mappa, disporre gli oggetti in
  pianta \\
\addlinespace
Left / Right (e le altre) & i lati, senza prospettiva, in \emph{Unlit} &
  controllare allineamenti, altezze, lati di una costruzione \\
\bottomrule
\end{tabular}
\caption{Le proiezioni della viewport viste a lezione.}
\label{gd14:tab:viste}
\end{table}

\subsection{Lit, Unlit, Wireframe: la profondità sono le ombre}

Le modalità di visualizzazione erano già state viste
(sezione~\ref{gd11:sec:viewmode}). Il docente ne riprende il significato con
un'osservazione che prepara la lezione:
\begin{itemize}
  \item \textbf{Lit} è tutto processato come versione finale, con tutti i calcoli
        della luce e gli effetti di post-processing;
  \item \textbf{Unlit} spegne l'illuminazione, e quindi non si vedono le ombre.
\end{itemize}
``L'illuminazione è un po' il contrario, in realtà, di quello che sta succedendo
dietro le quinte.'' \emph{Unlit} mostra gli oggetti ``al 100\% della loro
visualizzazione''; l'illuminazione aggiunge il calcolo delle \textbf{ombre}, e
\textbf{sono le ombre a creare la profondità}. Il docente si mette in una posizione
in cui un oggetto ne copre un altro: in \emph{Lit} si distinguono, in
\emph{Unlit} no, ``anche se è lì'', perché è la \textbf{differenza di colore tra
zone in ombra e zone in luce} che permette di percepire la tridimensionalità.

\paragraph{Wireframe e il terreno.} In \textbf{Wireframe} si vedono le mesh: la
\emph{sky sphere} (lo sfondo) e, dato che la mappa aperta è quella di default, il
\textbf{terreno}, che essendo la simulazione di un terreno realistico è molto più
denso di poligoni rispetto ai cubi. E si nota un'altra cosa: il terreno è
\textbf{gestito per livelli di dettaglio}. Più ci si avvicina a un blocco di
terreno, più è definito; più è lontano, meno triangoli ci sono. ``Vedete come si
trasforma?'' Oltre una certa distanza tutti quei dettagli non servono, e da una
visione normale la differenza non si nota neanche. Sono \textbf{trucchi di
ottimizzazione}.

\paragraph{Dell'ottimizzazione fregatevene.} Ed ecco una frase che il docente
ripeterà ``ripetutamente nel corso di quest'anno'':
\begin{quote}
Al primo anno \textbf{fregatevene totalmente dell'ottimizzazione}. Non fa parte dei
nostri problemi.
\end{quote}
Con ``qualche piccolissimo asterisco'', alcuni dei quali si vedranno proprio in
questa lezione (sezione~\ref{gd14:sec:troppe-luci}).

\subsection{Genitori, figli e componenti}

\paragraph{Il ripasso.} Ripasso veloce della gerarchia
(sezione~\ref{gd11:sec:gerarchie}): se si trascina un cubo più piccolo sopra un
cubo più grosso nell'Outliner, i due sono \textbf{imparentati}, e ogni
trasformazione del genitore, \textbf{scala inclusa}, viene ereditata dal figlio;
le modifiche al figlio, invece, non risalgono al genitore.

\paragraph{Figli e componenti non sono la stessa cosa.} C'è però una distinzione da
fare. Selezionando un oggetto nel mondo di gioco, nel pannello \emph{Details} si
vede l'\textbf{istanza}, e dentro l'istanza i suoi \textbf{componenti}. Un cubo ha un
solo componente, la mesh del cubo; un Blueprint più complesso, come un personaggio
o un oggetto che emette particelle, avrà sia la mesh sia l'emettitore di
particelle. ``Un actor è formato da più componenti.'' Quando si parla di
\textbf{genitore e figlio}, invece, non si parla di componenti, ma di
\textbf{actor}.

\paragraph{La domanda sulla PlayStation.} Federica chiede: nel fare la PlayStation
ha costruito il modello pezzo per pezzo, imparentando i pezzi; sarebbe stato
meglio farne un componente unico, pensando alla programmazione? La risposta: tutti
questi esercizi servono a mettere le mani, e ``nessuno dei metodi che avete usato
finora è un modo che si usa per fare un videogioco''. Quella PlayStation
``l'avrebbe fatta un modellatore 3D in Maya o in Blender''. Quello che contava era
far passare l'idea dell'\textbf{ereditarietà}: il figlio eredita tutti i cambiamenti
del genitore. ``Non preoccupiamoci troppo dei dettagli per ora; impariamo a
conoscere le cose con cui stiamo lavorando.''

\paragraph{Selezionare la discendenza.} Un'ultima comodità: facendo tasto destro su
un genitore, \textbf{Select All Descendants} seleziona tutti i discendenti,
mentre \textbf{Immediate Children} seleziona solo i figli diretti e non ``i figli
dei figli''. Serve per selezionare velocemente più oggetti insieme (sulla slide
sono indicate anche le relative scorciatoie da tastiera\footnote{Nella trascrizione
``Ctrl+Shift+D e Ctrl+Alt+D'', senza che sia chiaro a quale comando corrisponda
ciascuna: si rimanda alla slide.}).

% ---------------------------------------------------------------------
\section{Il Physically Based Rendering}
\label{gd14:sec:pbr}

``Parliamo ora di una cosa, e questa sarà una sezione un po' teorica, ma è
importante capire cosa sta succedendo nell'engine'': il \textbf{PBR},
\emph{Physically Based Rendering}.\footnote{Nella trascrizione anche ``PDR''.}

\subsection{Molti modi per disegnare un oggetto}

Quando si parla di grafica 3D, ci sono molti modi per rappresentare un oggetto
tridimensionale su schermo, e nella storia dei videogiochi si sono usate tecnologie
molto diverse per raggiungere lo stesso obiettivo, anche per il 2D. Chi ha giocato
a un gioco 2D della prima PlayStation, come \emph{Castlevania: Symphony of the
Night}, ha in realtà giocato a un \textbf{gioco 3D} in cui, invece di modelli
tridimensionali, ci sono \textbf{quadrati con sopra gli sprite}. È il motivo per cui
dentro quei giochi compaiono piccoli elementi 3D che ruotano, come la stanza di
salvataggio con ``quella specie di dado che gira''.

\paragraph{La definizione.} Il metodo standard moderno è il \textbf{Physically Based
Rendering}. Con \emph{rendering} si intende la realizzazione di un'immagine a
schermo; il PBR è un metodo di rendering che si usa per \textbf{approssimare
l'interazione realistica della luce con gli oggetti}.

\paragraph{Un problema di ordine.} Il PBR è un modo di applicare la luce che
dipende dai \textbf{materiali}, e la lezione sui materiali non c'è ancora stata.
``Questa è una cosa che purtroppo tornerà spesso'': i videogiochi sono ``una roba
molto olistica'', da considerare in ogni componente, ``e io non posso farvi una
sola lezione di sedici ore che copre tutto quanto''. Capiterà quindi che, per
introdurre un concetto, si debba toccare un concetto di cui non si è ancora
parlato: il docente cercherà di farne accenni ``il più comprensibili possibile''.

\subsection{Texture e materiale, di nuovo}

``Antonino, che cos'è una texture?'' La risposta dello studente (la superficie
dell'oggetto) è ``esattamente quello che volevo tirar fuori, perché non è quello,
una texture. Ma una volta era così.''

\paragraph{Una volta.} Una volta la texture era, ``detto terra-terra'',
\textbf{quello che riempie lo spazio fra i triangoli}. Una mesh è fatta di
triangoli (sezione~\ref{gd11:sec:mesh}); in modalità \emph{Wireframe}, togliendo le
icone con il tasto G, si vedono i triangoli che compongono le mesh. L'idea era che
la texture fosse la superficie proiettata su questi triangoli.

\paragraph{Oggi.} Oggi la texture è \textbf{solo una delle componenti} di quella
superficie. Una texture è un'\textbf{immagine}: un file quadrato, di lato pari a
una potenza di 2 (``ma non entriamo nei dettagli tecnici''), in PNG o in altri
formati, che viene preso dal \textbf{materiale}, ``passato attraverso tutta una
serie di cose'' e applicato all'oggetto. Spesso \emph{texture} e \emph{materiale}
sono usati come sinonimi, ``e questa confusione è normalissima, tutti ce l'hanno''.
Ma la distinzione è importante:
\begin{quote}
\textbf{Un materiale usa texture.} E nel PBR il materiale divide ciò che viene
applicato ai triangoli in \textbf{più parti}, che calcolate nel loro insieme danno
il risultato.
\end{quote}

\subsection{Che cosa racconta la luce}

\paragraph{La sedia.} Il progetto aperto ha lo Starter Content, quindi la
dimostrazione si fa sulla ``nostra ben amata sedia, che vedremo molto spesso''
(tasto \textbf{F} per inquadrare l'oggetto selezionato). ``Guardate in alto a destra
sulla sedia: cosa sta succedendo con la luce?'' Muovendo la telecamera, i
riflessi chiari \textbf{si spostano}. Come mai? Sono creati dal riflesso del sole,
che forma un angolo con la posizione di chi osserva. Tecnicamente si chiamano
\textbf{specular highlights} (``ma ci interessa relativamente'').

\paragraph{Informazioni fisiche.} Sono \textbf{informazioni fisiche}, ed è per questo
che si parla di rendering \emph{basato sulla fisica}: la luce ci sta passando
informazioni sulla fisica dell'oggetto, e quella fisica dice ``questa è una
superficie riflettente''. Facendo la stessa cosa con una \textbf{roccia} (``grandioso,
sono un genio, l'ho messa in ombra''), c'è comunque un certo livello di interazione
con la luce, ma \textbf{non ci sono specular highlights}, nessun punto di riflesso
del sole. Perché sono due materiali fisicamente diversi: la roccia è riflettente
solo fino a un certo punto, la plastica e il metallo sono molto riflettenti, cioè
\textbf{lucidi}.

\paragraph{Il vetro e il metallo.} Il PBR permette di passare a un oggetto le
informazioni che il calcolo della luce userà per \textbf{simulare il materiale di cui
è fatto}. Per il \textbf{vetro} dice anche altro: che è trasparente e che, grazie a
uno \emph{shader} (``lì andiamo in un campo un po' più complesso''), distorce la
visuale di quello che c'è dietro. La base della sedia, invece, è di
\textbf{metallo}, con la forte capacità di riflettere la luce tipica del metallo.
\begin{quote}
Il Physically Based Rendering è \textbf{un sistema unitario di ricreazione delle
caratteristiche fisiche naturali delle superfici}.
\end{quote}
Nella slide c'è un link alla documentazione di Unreal per chi vuole approfondire.
La parte importante da ricordare è che \textbf{un materiale è composto da più
elementi calcolati insieme} (figura~\ref{gd14:fig:pbr}).

\subsection{L'albedo: solo il colore}

La base è l'\textbf{albedo}, detto anche \emph{color}, \emph{base color} o
\emph{diffuse}. ``Che cos'è l'albedo? È il colore, \textbf{unicamente il
colore}.'' I nomi sono tanti perché i termini sono sempre la stessa cosa ma, a
seconda del programma, si chiamano diversamente; più spesso \emph{albedo} o
\emph{color}. (In ambito scientifico \emph{albedo} indica altro, per esempio la
capacità riflettente dei pianeti; nel nostro settore è la texture che rappresenta
unicamente i colori di un oggetto.)

\paragraph{Unlit è l'albedo.} Sulla slide, a destra, il materiale di un oggetto
metallico è renderizzato mostrando solo l'albedo. ``Questo è quello che succede
quando fate \emph{Unlit}'': si vede unicamente il colore, senza nessun calcolo di
come la luce impatta le superfici in base al materiale di cui sono fatte.

\paragraph{Dentro il materiale.} Il docente apre velocemente il materiale della
roccia (doppio clic sul materiale, non sull'oggetto). Il materiale è composto da
più nodi; uno di questi è il \emph{Texture Sample} collegato a \emph{Base Color}:
è la texture dell'albedo, chiamata \texttt{T\_Rock\_Basalt\_D}. Cliccando sulla
piccola icona accanto, il Content Browser si apre direttamente sulla texture. La
\texttt{D} finale sta per \emph{diffuse}: di nuovo, ``la texture dei colori di un
oggetto''.
\begin{itemize}
  \item ``Lo so, è un po' antipatico lavorare in un settore in cui una cosa può
        avere tre nomi diversi'', ma ci si abitua, e soprattutto, quando la si
        vede, si capisce subito qual è la texture del colore.
  \item Nell'editor della texture ci sono i canali \textbf{R}, \textbf{G},
        \textbf{B} (rosso, verde, blu) e l'\textbf{alfa}, la trasparenza, che il
        docente disattiva perché per ora non interessa. Lasciando acceso un solo
        canale, Unreal mostra l'immagine \textbf{in bianco e nero}, perché è più
        facile da leggere.
\end{itemize}
``Non preoccupiamoci: le vedremo nella lezione dedicata a materiali e texture.
Queste prime lezioni sono uno scarico enorme di roba da sapere, lo so, me ne rendo
conto: se ci sono dubbi, interrompetemi.''

\subsection{La roughness: quanto è ruvida una superficie}

Un'altra parte del PBR è la \textbf{roughness}\footnote{Nella trascrizione anche
``Raphness''.}: quanto è \textbf{ruvida} una superficie. Sulla slide, in basso, lo
stesso oggetto ha un colore di base più la roughness, con valori crescenti: più
aumenta la roughness, \textbf{meno l'oggetto riesce a riflettere}.

\paragraph{Valori, non categorie.} Il PBR non funziona con ``questo materiale
plastica, questo metallo, questo roccia, questo carta''. Ha \textbf{una serie di
valori}, e in base alla quantità di ciascuno si approssima un tipo di superficie.
Il docente lo fa sperimentare: ``Francesco, che materiale è questo primo oggetto?''
Acciaio, un qualche metallo. E questo? Pietra, o plastica. E questo? ``Io direi
plastica, metallo opaco, qualcosa del genere.'' Cambiando \textbf{un solo valore}
si stabilisce che tipo di materiale si sta creando.

\paragraph{Un numero o una texture.} Il valore può essere \textbf{un solo numero}
applicato a tutto l'oggetto, oppure può essere preso da una \textbf{texture} che
dice ``queste parti sono di metallo, queste di plastica'': una texture in
\textbf{scala di grigi} che dice quanta roughness applicare in ciascun punto. Per
questo la distinzione tra texture e materiale è importante:
\begin{quote}
Un materiale può essere composto da più texture, da una sola texture e dei numeri,
o solo da numeri. Un materiale è \textbf{una scatola} in cui mettiamo quello che
vogliamo vedere su schermo.
\end{quote}

\paragraph{Perché la ruvidezza spegne i riflessi.} Che cosa fa la roughness per
ottenere questo effetto? L'illustrazione della slide lo mostra: invece di far
rimbalzare la luce in modo uniforme, \textbf{più una superficie è ruvida, più fa
rimbalzare la luce in direzioni casuali}. È una superficie piena di imperfezioni,
tanto che toccandola si sente che è ruvida, perché si sentono le imperfezioni. Una
plastica ruvida non fa riflessi, una plastica lucida sì; lo stesso vale per i
pavimenti e per un sacco di altre cose.

\subsection{Lo specular}

C'è poi lo \textbf{specular}, dal latino \emph{speculum}, ``che vuol dire
specchio'': indica \textbf{quanto è riflettente} un materiale. I riflessi sono un
discorso particolare, e roughness e specular \textbf{si influenzano a vicenda}: se
c'è troppa roughness, lo specular non funziona più. ``Non vi serve saperlo per
adesso'': è solo un'introduzione per capire come funziona la luce in Unreal,
perché tutti questi valori \textbf{toccano la luce} che si dà nell'illuminazione.

\paragraph{La domanda: che differenza c'è?} Uno studente chiede: visivamente, la
differenza tra la ruvidezza e lo specular qual è? Sembrano due modi per ottenere
la stessa cosa, e allora perché usare uno e non l'altro? ``La risposta breve è che
\textbf{sono quasi uguali}, sì. La risposta lunga sono questioni di
\emph{technical art}'' che sarebbero una parentesi troppo lunga. Ci saranno due
lezioni dedicate ai materiali: ``questa domanda te la segno con un asterisco e la
rivediamo quando andiamo a vedere i materiali''.

\subsection{La mappa della roughness e i bordi rovinati}

Nell'immagine a destra della slide, con le texture delle varie caratteristiche
applicate, la \textbf{mappa della roughness} dice che l'oggetto è ruvido
\textbf{solo sugli spigoli}, per dare l'effetto di bordi un po' rovinati. ``Se date
un'occhiata alla chiavetta USB che vi ho appena fatto girare, forse avrete notato
che ha i bordi rovinati'': serve a dare la sensazione di \textbf{cosa vissuta}.
Guardando da vicino, la superficie è riflettente dove c'è il metallo o la vernice,
ma sui bordi è più ruvida, e la mappa della roughness segna proprio quei bordi.

\begin{figure}[htbp]
\centering
\begin{tikzpicture}[
  gd14lab/.style={font=\scriptsize, align=center},
  gd14box/.style={draw, rounded corners=2pt, minimum width=2.4cm,
                  minimum height=0.8cm, font=\scriptsize\bfseries, align=center},
]
  % riga delle sfere: roughness crescente
  \foreach \i/\r in {0/0, 1/0.25, 2/0.5, 3/0.75, 4/1} {
    \pgfmathsetmacro\hl{100-90*\r}
    \pgfmathsetmacro\hs{0.32-0.18*\r}
    \shade[ball color=gray!60] (\i*1.9,0) circle (0.6);
    \fill[white, opacity=\hl/100] (\i*1.9-0.22,0.22) circle (\hs);
    \node[gd14lab] at (\i*1.9,-0.95) {roughness\\$\r$};
  }
  \draw[-{Stealth}] (-0.6,-1.55) -- (8.2,-1.55)
    node[midway, below, font=\scriptsize] {più ruvida $\Rightarrow$ meno
    riflessi, specular highlight sempre più larghi e deboli};
  % materiale come scatola
  \node[gd14box, fill=orange!20] (alb) at (0,-3.2) {Albedo\\(colore)};
  \node[gd14box, fill=gray!25] (rou) at (2.9,-3.2) {Roughness\\(ruvidezza)};
  \node[gd14box, fill=cyan!15] (spe) at (5.8,-3.2) {Specular\\(riflettenza)};
  \node[gd14box, fill=gray!10, minimum width=2.0cm] (alt) at (8.4,-3.2)
     {altri canali\\(normali, \ldots)};
  \node[draw, very thick, rounded corners=4pt, fit=(alb)(rou)(spe)(alt),
        inner sep=6pt, label={[font=\scriptsize\bfseries]above:Materiale: una
        scatola di texture e numeri}] (mat) {};
  \node[gd14box, fill=yellow!25, minimum width=3.2cm] (luce) at (4.2,-5.1)
     {calcolo della luce\\(il risultato a schermo)};
  \draw[-{Stealth}, thick] (mat.south) -- (luce.north);
\end{tikzpicture}
\caption{Il PBR in sintesi. In alto: lo stesso colore con roughness crescente;
cambiando un solo valore si passa da un metallo lucido a una superficie opaca. In
basso: il materiale come ``scatola'' di più elementi (ciascuno un numero o una
texture) che il calcolo della luce usa insieme. Ricostruzione schematica.}
\label{gd14:fig:pbr}
\end{figure}

% ---------------------------------------------------------------------
\section{Le luci e la loro mobilità}
\label{gd14:sec:mobilita}

``Passiamo oltre. Luci.'' Delle luci si è già parlato la volta precedente, ``per
darvi un'infarinatura generale'' (sezione~\ref{gd11:sec:luci}). I tipi sono quattro,
più la \textbf{rect light} (luce rettangolare), che per ora si mette da parte:
\textbf{directional light}, \textbf{sky light}, \textbf{point light} e
\textbf{spot light}. Il docente decide di vederle non nel vuoto ma dentro
\emph{Medieval Village}, ``così possiamo vedere queste cose applicate in un contesto
già realizzato a un altissimo livello''. Mentre il progetto finisce di compilare
tutti i materiali (``dato che non l'ho mai aperto prima''), si fa una nota
importante.

\subsection{Static, Stationary, Movable}

Aprendo una luce, nei \emph{Details} c'è una sezione che dice se la luce è
\textbf{mobile} o no: la \textbf{Mobility}, con tre valori
(tabella~\ref{gd14:tab:mobilita}).

\paragraph{Static: il baking.} Le luci \textbf{statiche} sono calcolate \textbf{una
volta sola}: non si possono spostare durante il gameplay e non si possono spostare
nemmeno gli oggetti che illuminano, e creano delle \textbf{lightmap}. È una
distinzione che vale per moltissimi engine e che va sotto il nome di
\textbf{baking delle luci} (cottura). Dato che \textbf{calcolare la luce in tempo
reale è una delle cose più costose che fa un engine}, una delle cose che si
facevano agli inizi del 3D\footnote{Nella trascrizione ``all'inizio dell'epoca di
Jock 3D'', non ricostruibile con certezza.}, e che in molti settori si fa ancora,
era \textbf{precalcolare le luci}: le calcoli una volta e le converti in dati. Se
non serve spostare luci, aggiornare cose o spostare parti dello scenario, si può
dire: queste luci sono statiche, non cambieranno mai; calcolale una volta e non
cambiarle più. Unreal usa le \textbf{lightmap}: texture con tutta la luce del
livello, caricate e applicate ai vari oggetti.
\begin{itemize}
  \item \textbf{Vantaggio}: leggerissime, ``girano molto facilmente''.
  \item \textbf{Svantaggio}: non si può spostare niente. Se si sposta un oggetto da
        un punto buio a un punto illuminato, in un ambiente con luci statiche la
        sua illuminazione \textbf{non cambia}: le luci lo hanno illuminato nel
        momento in cui sono state calcolate. Spostandolo al sole, l'oggetto resta
        in ombra, ``perché c'è ancora solo il calcolo delle luci precedente''.
\end{itemize}

\paragraph{Lumen: tutto è dinamico.} Questa cosa, in Unreal, è di fatto
\textbf{disattivata}. Unreal usa un sistema di illuminazione dinamica che si chiama
\textbf{Lumen}\footnote{Nella trascrizione ``l'humen'', ``l'UMEN'', ``lume''.}, che
\textbf{considera automaticamente tutte le luci come dinamiche}. Il discorso sulla
mobilità, quindi, si applica davvero solo se Lumen è disattivato; il docente lo
racconta ``come informazione di servizio'' su cosa succede dietro le quinte,
perché potrebbe tornare utile, ma gli studenti di solito lavorano con Lumen
acceso.

\paragraph{Lo sfarfallio.} Se in una stanza si notano \textbf{sfarfallii di luce},
\textbf{rumore visivo} legato alle luci, o se in giochi moderni fatti con Unreal si
è notato uno strano sfarfallio ai confini tra luce e ombra, ``è Lumen, che calcola
l'illuminazione dinamica in tempo reale \textbf{approssimata}''. Detto in modo molto
semplice: invece di calcolarla perfetta a ogni frame, la approssima \textbf{tra un
frame e l'altro}, in modo da avere un'approssimazione del ray tracing senza dover
fare ray tracing (sezione~\ref{gd14:sec:lumen}).

\paragraph{Movable: la luce pura.} Una luce \textbf{Movable} è ``la luce pura'': si
può spostare la sorgente durante il gioco e l'illuminazione di ciò che la circonda
cambia; si possono spostare gli oggetti intorno e continueranno a proiettare ombre.
È la luce di base, senza limiti, che si comporta come una luce realistica, e come
si può immaginare è \textbf{la più costosa} in termini di risorse.

\paragraph{Stationary: a metà.} Una via di mezzo è la luce \textbf{Stationary}: una
volta piazzata nell'engine non può più essere mossa durante il gioco, ma continua
comunque a \textbf{proiettare le ombre} degli oggetti che le si spostano davanti.

\begin{table}[htbp]
\centering
\small
\begin{tabular}{@{}p{2.2cm}p{3.6cm}p{3.4cm}p{3.4cm}@{}}
\toprule
\textbf{Mobility} & \textbf{Che cosa fa} & \textbf{Pro} & \textbf{Contro} \\
\midrule
Static & calcolata una volta, salvata in lightmap (\emph{baking}) &
  leggerissima & niente si può muovere: un oggetto spostato al sole resta in
  ombra; va ``buildata'' a mano \\
\addlinespace
Stationary & la luce non si muove nel gioco, ma le ombre degli oggetti mobili sì &
  compromesso tra costo e libertà & la sorgente resta fissa \\
\addlinespace
Movable & luce pienamente dinamica: si muovono luce e oggetti &
  si comporta come una luce reale & la più costosa \\
\bottomrule
\end{tabular}
\caption{La mobilità delle luci in Unreal. Con \textbf{Lumen} attivo (il default,
e la scelta consigliata per il corso) tutte le luci sono trattate come dinamiche e
la distinzione conta solo in casi particolari.}
\label{gd14:tab:mobilita}
\end{table}

\paragraph{Sapere che esistono.} ``Finché usate Lumen, non preoccupatevi troppo di
queste cose.'' Quando si vorrà ottimizzare, o spingersi in situazioni più avanzate
di controllo delle luci, si potrà cominciare a guardarle. Ed ecco l'approccio del
docente a tutto il corso:
\begin{quote}
Io toccherò un sacco di cose in queste lezioni che, se dovessi spiegarle nei
dettagli, richiederebbero intere lezioni. Mi interessa di più farvi sapere che
queste cose \textbf{esistono}, in modo che, se avete bisogno di modificarle,
sappiate dove andare a cercare, cosa cercare su Google, dove sono le varie opzioni.
\end{quote}
``Sì, c'è una buona parte di \textbf{studio individuale} all'interno di queste
lezioni. È inevitabile.''

\paragraph{La scritta Preview e Build Lighting Only.} Se si lavora con Lumen spento
e sugli oggetti in scena, in mezzo alle ombre, compare la scritta
\textbf{Preview}, è perché Unreal sta mostrando un'anteprima delle luci statiche e
chiede, prima di giocarci, di \textbf{buildarle}, cioè di fare i calcoli di tutte
le luci adesso. La soluzione: menu \textbf{Build} $\to$ \textbf{Build Lighting
Only}. Con Lumen acceso questa scritta non comparirà ``praticamente mai''; ma chi
vuole provare la \emph{static lighting} per i fatti propri deve sapere che le luci
\textbf{vanno buildate a mano}.

\paragraph{Approfondimento: il baking secondo Totten.} In \emph{An Architectural
Approach to Level Design} Totten osserva che, pur dovendo tenere d'occhio il costo
grafico dell'illuminazione in tempo reale, un'illuminazione drammatica si può
ottenere con il \textbf{light baking}: il programma renderizza una mappa di colore
dell'ambiente che contiene le informazioni di illuminazione. Molti engine hanno
funzioni di \emph{light mapping} che producono una texture con le informazioni
sulla luce, riducendo la necessità di illuminazione in tempo reale nelle aree in
cui la luce non cambierà in modo significativo. È esattamente il compromesso della
tabella~\ref{gd14:tab:mobilita}: Lumen sposta l'ago della bilancia verso il
dinamico, ma il principio resta.

% ---------------------------------------------------------------------
\section{La directional light e il SunSky}
\label{gd14:sec:directional}

\subsection{Il sole: un muro di luce infinito}

``Francesca, che cos'è la directional light?''\footnote{Nella trascrizione
``directional line''; più avanti anche ``directionalite'' e ``detection light''.}
``La luce del Sole.'' Per essere tecnicamente più precisi, la directional light è
la luce legata al \textbf{ciclo diurno}: una fonte di luce in cui \textbf{i raggi
procedono tutti in parallelo}. ``Immaginatelo come un muro di luce dalle dimensioni
infinite'': è quello che è il Sole rispetto alla Terra. A livello pratico i raggi
solari sono paralleli, perché la differenza tra un raggio e l'altro è così minima
da essere trascurabile (figura~\ref{gd14:fig:tipi}).

\paragraph{Una sola.} ``Lo ripeto ancora una volta: \textbf{una sola directional
light} nel vostro gioco.'' Unreal non impedisce di metterne più di una, ma se lo si
fa ``finite per avere grossi problemi''. Assicurarsi che ce ne sia solo una.

\paragraph{Conta solo la rotazione.} L'unica cosa che si può fare è
\textbf{ruotarla}: essendo i raggi infiniti e paralleli, \textbf{spostarla non cambia
assolutamente nulla}. Ma proprio per questo la si può spostare a proprio comodo.
Se si sta lavorando a una casa, per esempio la casa del protagonista, che è il
punto di partenza per l'estetica di tutto il livello, e si vuole che la luce cada
bene su quella casa per poi costruire il resto in base a quella luce, ci si può
portare la directional light lì accanto, così da vedere nello stesso momento sia
la casa sia la luce mentre la si ruota.

\subsection{Il SunSky di Medieval Village}

Nel progetto di \emph{Medieval Village} la directional light è particolare. Tutte
le luci ambientali che normalmente stanno in una cartella (la cartella
\emph{Lighting}, sezione~\ref{gd11:sec:da-zero}) sono raccolte dentro un unico
oggetto chiamato \textbf{SunSky}: ``come un Blueprint che ha all'interno dei
figli'', tra cui la directional light. Ruotando la luce, \textbf{girano anche le
ombre} di tutta la scena; spostando l'oggetto SunSky, invece, non cambia nulla.
Una precisazione: qui la directional light è un \textbf{componente} del SunSky; nei
livelli degli studenti è un \textbf{actor} a sé (la distinzione della
sezione~\ref{gd14:sec:ripasso}).

\paragraph{Non ci capirete niente, ed è normale.} Il docente dà un primo sguardo al
villaggio: ``un ambiente medievale fatto giocabile'', a un livello di qualità ``da
produzione di altissimo livello''. Serve per due motivi: avere un punto di
riferimento, e poter guardare dietro le quinte come sono state fatte le cose. Ma
``se ci guardate dentro adesso, non ci capirete niente''. Innanzitutto il livello è
pieno di icone (si tolgono e si rimettono con \textbf{G}). E poi molte cose sono
fatte con strumenti che non si sono ancora visti: per esempio, cliccando sull'erba,
invece del \emph{landscape} si seleziona un pezzo di terreno. ``Non preoccupatevi
se non ci capite niente. È normale. A fine anno ci capirete molto di più.'' Si nota
anche la \textbf{nebbia} (\emph{fog}) in lontananza, ``che fa questo effetto così''.
``Tutto quello che volete studiare e con cui giocare è qui.''

\subsection{Il colore del sole, e la notte}

Nei \emph{Details} della directional light ci sono molti valori, e alcuni si sono
già visti. Il primo è il \textbf{colore}. ``Volete fare la scena finta notte? Lo
spostiamo un po' sul blu.'' E abbassando un po' la luce, la scena diventa notte:
\begin{quote}
Come a Hollywood, in realtà la notte non esiste come concetto. Che cos'è la notte?
\textbf{La notte è il giorno con una tinta blu.}
\end{quote}
Ovviamente, con un po' più di studio rispetto alla dimostrazione, ``in cui l'unica
cosa che ho fatto è cambiare la luce''. (È la tecnica cinematografica nota come
\emph{day for night}, la ``notte americana''.)

\subsection{Cascade Shadow Maps: i valori tecnici}

Tra i valori ce n'è uno che il docente mostra come esempio: nella sezione
\textbf{Cascaded Shadow Maps}, la \textbf{Dynamic Shadow Distance} indica la qualità
delle ombre: più è alto il valore, più le ombre sono \textbf{definite} a distanza.
Serve soprattutto a far vedere che i valori nei \emph{Details} non sono solo luce,
posizione e simili: ci sono anche molti valori tecnici. ``Non andate a toccarli
per ora.''

\paragraph{Uno screenshot a un frame al secondo.} ``Tutto quello che stiamo facendo
noi è fare screenshot. Se anche il vostro computer va a un frame al secondo, basta
per fare uno screenshot.'' Quindi, per il primo anno, niente preoccupazioni di
ottimizzazione.

\paragraph{Quando un valore non fa niente.} Una cosa da ricordare sempre: se si
tocca un valore di una luce e non cambia niente, c'è la possibilità che quel valore
influenzi \textbf{solo certi tipi di luce} (solo le statiche, solo le stazionarie e
le mobili, e così via). Unreal non mostra un pop-up che avvisa ``guarda che questo
valore tocca solo le luci mobili''. Però, \textbf{lasciando il mouse sopra} un
valore, compare spesso una descrizione di cosa fa e, nel caso delle luci, per
quale mobilità vale. ``Tenetele d'occhio, queste cose, perché tendono a essere le
robe in cui modifichi e dici: non sta succedendo niente, come mai?''

% ---------------------------------------------------------------------
\section{La Sky Light e il colore delle ombre}
\label{gd14:sec:skylight}

``Vi ricordate che l'avevo accennata la scorsa volta? Che cos'è la Sky
Light?''\footnote{Nella trascrizione ``Skylite''.} Daniele risponde: dà i colori
della luce nell'atmosfera. La \textbf{Sky Light} è la \textbf{diffusione della luce
nell'atmosfera}: è il motivo per cui il cielo è azzurro invece di essere nero,
perché la luce si diffonde nell'atmosfera e, diffondendosi, dà
un'\textbf{illuminazione di fondo generale}: il tono generale del colore della luce
che si riflette nel cielo.

\paragraph{Detto in modo brutale: le ombre.} ``In una maniera un po' più brutale, la
Sky Light sono le ombre.'' Cioè: è \textbf{la luce dove non arriva la luce diretta del
sole}. Nella lezione precedente, costruendo un livello da zero, si era visto che
con il solo sole (tolto il discorso dell'illuminazione dinamica approssimata di
Lumen) un oggetto è illuminato da un lato solo
(sezione~\ref{gd11:sec:da-zero}): è il motivo per cui la Luna ha lo spicchio, perché
il sole non la illumina dall'altro lato, che quindi è totalmente nero. La Sky
Light è responsabile di \textbf{che cosa c'è dentro le ombre} (``in parte, perché
tutti i sistemi si sovrappongono'').

\paragraph{Cambiare il colore delle ombre.} Questo significa che, se per motivi
stilistici si vuole cambiare il \textbf{colore}, la ``temperatura'', delle ombre, si
seleziona la Sky Light e si cambia il suo colore. Il docente lo mostra in
\emph{Medieval Village}:
\begin{itemize}
  \item un \textbf{horror medievale}, con un'atmosfera realistica ma più cupa: Sky
        Light verso il \textbf{rosso}. Un rosso acceso ``è veramente pesante'';
        abbassato un po', dà alla scena un tono leggermente rosso, ``un po' più
        tetro'';
  \item un villaggio coperto dal ``\textbf{muschio elfico}'': Sky Light sul
        \textbf{verde}.
\end{itemize}
Spinti al massimo sono ``un pugno in un occhio'', ma servono a far capire come
funziona. Poi si riporta la Sky Light al \textbf{bianco}, la luce neutra. (Il
risultato è attenuato dal fatto che il post-processing di questo livello è molto
\textbf{desaturato}, e toglie parecchio colore in generale: ``ma ce ne
preoccuperemo più avanti''.)

\paragraph{Approfondimento: un'unica luce, molte atmosfere.} Totten racconta il
caso del gioco per dispositivi mobili \emph{SWARM!}, in cui i limiti dei dispositivi
permettevano \textbf{una sola directional light} per livello. I designer dovevano
quindi piazzarla con cura, all'angolo da cui volevano che arrivasse la luce, e
regolare il \textbf{colore della luce ambientale}, l'intensità e il colore della
luce stessa: con qualche ritocco creativo si ottenevano condizioni di luce molto
diverse con una sola luce. È la stessa lezione della dimostrazione: rotazione e
colore del sole più il colore della luce ambientale bastano a cambiare l'ora del
giorno e il genere del gioco.

% ---------------------------------------------------------------------
\section{La point light e la verosimiglianza: tre esempi da Lies of P}
\label{gd14:sec:liesofp}

\subsection{La point light, di nuovo}

``Ancora un po' di luce. Point light, l'abbiamo già vista. Cos'è una point light?''
Una luce \textbf{radiale}: un punto specifico che emana luce \textbf{in tutte le
direzioni} attorno a sé (figura~\ref{gd14:fig:tipi}), e che proietta ombre tutto
intorno.

\paragraph{F11.} ``Pro tip di Unreal, e seguitemi perché questo è utilissimo'':
premendo \textbf{F11} l'editor va \textbf{a schermo intero}, così si posizionano le
cose più comodamente. In questo caso Unreal avvisa anche che renderizzare a 4K una
scena del genere significa ``\emph{expect extremely poor performance}'': infatti
il computer va a pochi fps. (Della \textbf{scalabilità} dell'engine si parlerà in un
altro momento; per ora, di nuovo, niente ottimizzazione.)

\begin{figure}[htbp]
\centering
\begin{tikzpicture}[
  gd14ray/.style={-{Stealth[length=4pt]}, thick, orange!80!black},
  gd14cap/.style={font=\scriptsize\bfseries, align=center},
]
  % directional
  \begin{scope}
    \foreach \y in {0,0.45,...,2.25} {
      \draw[gd14ray] (0,\y) -- (2.4,\y-0.9);
    }
    \fill[gray!40] (0.2,-1.3) rectangle (2.8,-1.1);
    \node[gd14cap] at (1.4,-1.8) {Directional\\raggi paralleli, infiniti:\\conta solo la rotazione};
  \end{scope}
  % point
  \begin{scope}[xshift=4.6cm]
    \fill[yellow!80!orange] (1.2,0.8) circle (0.12);
    \foreach \a in {0,30,...,330} {
      \draw[gd14ray] (1.2,0.8) ++(\a:0.2) -- ++(\a:1.0);
    }
    \node[gd14cap] at (1.2,-1.8) {Point\\da un punto, in ogni\\direzione (una sfera)};
  \end{scope}
  % spot
  \begin{scope}[xshift=8.8cm]
    \fill[yellow!80!orange] (1.2,1.8) circle (0.12);
    \fill[yellow!25] (1.2,1.8) -- (0.1,-1.0) -- (2.3,-1.0) -- cycle;
    \fill[yellow!50] (1.2,1.8) -- (0.6,-1.0) -- (1.8,-1.0) -- cycle;
    \foreach \x in {0.3,0.75,1.2,1.65,2.1} {
      \draw[gd14ray] (1.2,1.7) -- (\x,-0.9);
    }
    \node[gd14cap] at (1.2,-1.8) {Spot\\da un punto, dentro un cono\\(cono interno ed esterno)};
  \end{scope}
\end{tikzpicture}
\caption{Tre delle luci viste a lezione. La Sky Light non ha una forma: è la luce
diffusa dall'atmosfera che riempie le zone d'ombra. Ricostruzione schematica.}
\label{gd14:fig:tipi}
\end{figure}

\subsection{Il barile che brucia}

Sulla slide c'è uno screenshot. ``Chi conosce questo gioco?'' È \emph{Lies of
P}.\footnote{Nella trascrizione ``IsoFP'', ``Lise of P'', ``Eyes of Pie'', ``Light
of P'', ``Light of Python''.} In una strada c'è un \textbf{barile} con dentro il
fuoco. ``Parliamo solo di luce, non di materiali: cosa sta succedendo qui?'' Gli
studenti partono a dire come lo avrebbero fatto loro; il docente li ferma: ``non vi
sto chiedendo come l'avreste fatto, vi sto chiedendo che cosa c'è in questa
immagine''. C'è una point light sul fuoco, certo. Ma \textbf{dove}?

\paragraph{Dove l'hanno messa.} Guardando le ombre proiettate, la point light è di
sicuro lungo una certa linea, e precisamente \textbf{sopra il barile}, non dentro.
``Non è un po' strano, se ci riflettiamo?'' A livello di verosimiglianza dello
scenario, la luce dovrebbe venire dall'origine del fuoco, cioè da dentro il barile.

\paragraph{Perché verosimiglianza e non realismo.} ``E qui iniziamo a vedere perché
io preferisco parlare di \textbf{verosimiglianza} e non di \textbf{realismo}.'' Nella
realtà:
\begin{enumerate}
  \item un barile che brucia fa \textbf{pochissima luce} (``ma quello è un discorso
        a parte'');
  \item se la sorgente fosse \textbf{dentro} il barile, tutta la parte sotto
        sarebbe in ombra, perché la luce sarebbe schermata dal barile stesso. Nella
        foto non si distinguerebbe nemmeno il barile, si vedrebbe solo il fuoco.
\end{enumerate}
In \emph{Lies of P} la luce è stata messa sopra perché \textbf{lo scopo del barile è
fornire una giustificazione per l'illuminazione della zona}. Messa dentro, la
luce renderebbe invisibile ciò che conta; messa sopra, giustifica una point light
fortissima che illumina tutta quella sezione della città. E se si mettesse un
barile qui e uno là, con la luce dentro? Si proietterebbero ``un sacco di buchi''
d'ombra lungo le pareti, ``bruttissimo da vedere''. La point light è piazzata in un
modo che illumina l'area \textbf{e} dà l'impressione che la luce venga dal barile,
anche se è fisicamente impossibile (figura~\ref{gd14:fig:barile}).

\begin{quote}
Non impuntatevi sul realismo, ragionate sulla verosimiglianza. \textbf{Non stiamo
facendo fotografie, stiamo facendo film.} I film illuminano con uno scopo
artistico. Il nostro obiettivo non è il realismo, è la verosimiglianza.
\end{quote}

\begin{figure}[htbp]
\centering
\begin{tikzpicture}[
  gd14cap/.style={font=\scriptsize, align=center, text width=5.2cm},
]
  % --- realistico ---
  \begin{scope}
    \fill[gray!65] (-0.2,0) rectangle (5.0,3.2);
    \fill[yellow!35, opacity=0.7] (2.4,1.0) -- (1.2,3.2) -- (3.6,3.2) -- cycle;
    \draw[thick] (-0.2,0) -- (5.0,0);
    \fill[gray!40] (2.0,0) rectangle (2.8,1.0);
    \fill[yellow!80!orange] (2.4,0.85) circle (0.1);
    \node[font=\scriptsize, text=white] at (0.9,0.4) {buio};
    \node[font=\scriptsize, text=white] at (3.9,0.4) {buio};
    \foreach \a in {60,75,90,105,120} {
      \draw[-{Stealth}, orange!80!black] (2.4,0.95) -- ++(\a:1.4);
    }
    \node[gd14cap] at (2.4,-0.7) {\textbf{Realistico}: la luce dentro il barile.
      Il barile scherma la luce: la strada resta al buio, si vede solo il
      fuoco.};
  \end{scope}
  % --- verosimile ---
  \begin{scope}[xshift=6.4cm]
    \fill[gray!15] (-0.2,0) rectangle (5.0,3.2);
    \begin{scope}
      \clip (-0.2,0) rectangle (5.0,3.2);
      \fill[yellow!30, opacity=0.8] (2.4,2.1) circle (2.3);
    \end{scope}
    \draw[thick] (-0.2,0) -- (5.0,0);
    \fill[gray!70] (2.0,0) rectangle (2.8,1.0);
    \fill[orange!70] (2.4,1.05) ellipse (0.35 and 0.12);
    \fill[yellow!80!orange] (2.4,2.1) circle (0.12);
    \node[font=\scriptsize\bfseries] at (3.3,2.4) {point light};
    \draw[dashed] (2.4,2.0) -- (2.4,1.1);
    \node[gd14cap] at (2.4,-0.7) {\textbf{Verosimile} (\emph{Lies of P}): la point
      light è sopra il barile. La strada è illuminata, e il fuoco ``giustifica''
      la luce.};
  \end{scope}
\end{tikzpicture}
\caption{Il barile di \emph{Lies of P}: dove starebbe la luce nella realtà e dove la
mette il level designer. Ricostruzione schematica.}
\label{gd14:fig:barile}
\end{figure}

\subsection{La lampada senza ombra}

Secondo esempio, sempre da \emph{Lies of P}: una lampada a muro. ``Da dove viene la
luce, a livello di realismo?'' Da dentro la lampada. ``E c'è qualcosa che manca?''
Manca \textbf{l'ombra}. Se la luce venisse davvero da dentro la lampada, dov'è
l'ombra della lampada stessa proiettata sul muro? Invece la luce si proietta ad arco
``come se la lampada non esistesse''.

\paragraph{Non proiettare ombre.} È una cosa che Unreal permette: si possono segnare
degli oggetti come \textbf{non proiettanti ombre}, ed è uno strumento usato molto
spesso dai level designer per evitare ombre brutte che non si vuole che compaiano.
``Nessuno si fermerà a dire: che schifo, non ci sono le ombre realistiche della
fonte di luce.'' Molto meglio così che avere ombre che si proiettano a caso lungo
le pareti e illuminano tutto male. ``Perché? Perché noi stiamo simulando quello
che sarebbe il \textbf{riflettore fuori dallo schermo} di un'illuminazione da
cinema.'' Non un'illuminazione iperrealistica, ma un'illuminazione verosimile.
Infatti, lì accanto, l'insegna del negozio (un'aragosta rossa\footnote{Nella
trascrizione ``l'insegna del Red Lobster'': si riporta come descrizione
dell'insegna.}) \textbf{l'ombra ce l'ha}: la scelta è fatta caso per caso.

\subsection{La luce motivata del corridoio}

Ultimo esempio prima della pausa: un corridoio con una lampada. ``Ignorate il
riflesso della luce sul personaggio protagonista. C'è qualcosa di particolare?''
Uno studente nota che la luce sembra arrivare da un altro punto. È una \textbf{luce
motivata} (``ma lo vedremo in un'altra lezione''): nella finzione del gioco, l'area
deve essere illuminata per \textbf{necessità di gameplay}, perché il gameplay si deve
vedere. Il level designer ha messo una lampada per dire ``qui c'è una luce'': il
giocatore vede la lampada e pensa ``ah, ok, la luce viene da quella lampada''.

\paragraph{Dov'è la point light?} Ma la point light è \textbf{altrove}: dove c'è
l'ombra, cioè più o meno al centro del corridoio. Si crea così un effetto curioso:
passando davanti alla lampada, l'ombra del personaggio si proietta \textbf{davanti}
alla lampada. ``È brutto? Sì.'' In \emph{Lies of P} ci sono un po' di punti in cui,
con un po' di esperienza di level design, si capisce che il passaggio di
sistemazione delle luci è stato fatto di fretta; ma l'esempio è ottimo per far
vedere che queste cose \textbf{si fanno}.

\paragraph{Non dentro la lampada.} ``Quello che voglio che vi entri in testa'' è che
non bisogna fare ``la roba del: ogni volta che ho una lampada, ci metto la luce
dentro, problema risolto''. Altrimenti si finisce per chiedere al docente perché
tutti i propri livelli hanno luci orribili.
\begin{quote}
Andate a vedere come si illumina un film: al di fuori della telecamera ci sono
cinquanta luci diverse che illuminano le cose che servono. Nei videogiochi è la
stessa cosa: le luci non si usano solo dove c'è la lampada, si usano per
\textbf{illuminare in base alle necessità estetiche della scena}.
\end{quote}
In \emph{Lies of P} si è fatta spesso questa scelta: sacrificare un approccio
realistico per un approccio verosimile. Lo scopo è che il corridoio sia ben
illuminato. ``Questi dettagli li notano pochissimi giocatori, quasi nessuno ci fa
caso. Ma intanto la composizione delle mie luci è stata fatta in una buona
maniera.''

\paragraph{Approfondimento: la luce che guida.} Totten cita l'architetto
Christopher Alexander, che in \emph{A Pattern Language} descrive come la luce possa
essere usata \textbf{agli incroci} dove il progettista vuole condurre le persone:
la luce come dispositivo per ``tirare'' gli occupanti attraverso lo spazio. E in un
altro passaggio del libro si ricorda che siamo attratti dagli elementi che
risaltano visivamente, che sembrano più importanti per scala, illuminazione o
contenuto: il level designer deve mantenere una \textbf{gerarchia di importanza},
attenuando le parti del livello che non sono primarie. È lo stesso ragionamento del
corridoio di \emph{Lies of P}: la luce non è dove sarebbe ``davvero'', ma dove serve
al giocatore.

\paragraph{La pausa.} ``Ok, pausa: ci diamo quindici, venti minuti. Se avete voglia di
farvi un caffè, fatevi un caffè.''

% ---------------------------------------------------------------------
\section{La pausa: ancora GTA, e una domanda sugli oggetti piazzati via codice}
\label{gd14:sec:pausa}

Durante la pausa la conversazione con alcuni studenti torna su \emph{GTA}. Anche
qui il capitolo riporta ciò che riguarda il design e l'industria.

\subsection{Monetizzazione e piedistalli}

\paragraph{Gli abbonamenti.} Uno studente, che stava leggendo le notizie, riferisce
che con l'abbonamento a \emph{GTA} (circa tredici dollari al mese) si ottengono
oggetti anche nel single player, e che per avere ``il gioco vero'' bisogna pagare
cento euro, altrimenti mancano missioni e certi negozi restano chiusi. Il docente
conferma che la serie è ``monetizzata in maniera folle'', e che l'online sarà
ancora peggio: in \emph{GTA Online} i soldi del gioco si possono ottenere giocando,
ma si possono anche \textbf{comprare} (le \emph{Shark Card}).

\paragraph{Tutti vogliono che sia un successo.} ``Qui ti sto dicendo una mia
opinione.'' Uno dei problemi dell'industria è che \textbf{molte volte si cerca un
piedistallo} su cui mettere qualcosa. Tutti vogliono che \emph{GTA 6} sia un
successo, per una lunga serie di motivi:
\begin{itemize}
  \item i \textbf{produttori}, perché sperano che sia un ritorno alle condizioni
        economiche di cinque--dieci anni fa: il modello AAA ``sta scricchiolando
        sempre di più, se non è già crollato'', il modello del singolo gioco che
        vende un sacco di copie e per cui tutti comprano una PlayStation;
  \item i \textbf{media}, che sperano in un ritorno all'epoca in cui ci si
        concentrava sui grandi giochi singoli, invece che su quelli come i giochi
        di Hoyoverse\footnote{Nella trascrizione ``i giochi di Mikoyo, ai
        Gaia'': probabilmente HoYoverse (miHoYo); non ricostruibile con
        certezza.} o i giochi per cellulare che stanno prendendo piede;
  \item i \textbf{giocatori casual}, perché \emph{GTA} è ``il buon nome''.
\end{itemize}
Questo potrebbe portare ad avere ``un po' di paraocchi'' e a ignorare certe cose.

\paragraph{La varietà, non la profondità.} Resta il fatto che, a livello di
contenuti, di quantità, di tipi di gameplay, di varietà, di sistemi (centinaia di
veicoli diversi), ``lo può fare solo Rockstar''. In quanti altri giochi, se si ha
voglia di fare il vigilante, si può rubare un'auto della polizia e fare le missioni
di arresto, o prendere un camion dei pompieri e spegnere gli incendi, ricevendo
soldi da spendere nella trama principale? È una serie con \textbf{poca profondità
lato meccaniche, ma una varietà allucinante}. Non esistono i montaggi di trickshot
o di sniping come in \emph{Call of Duty}: tutte le meccaniche di tutti i \emph{GTA}
sono sempre state \textbf{funzionali allo spazio di gioco}, non alla sfida meccanica
in sé. E questo aiuta il pubblico casual: non ci si deve allenare per giocare a
\emph{GTA} come ci si allena per \emph{Counter-Strike} o \emph{Call of Duty}. Ci si
gioca dall'inizio alla fine facendo un po' quello che si vuole: ``vuoi fare il
tassista? Vai a fare il tassista.''

\paragraph{Un finto MMO.} ``Per questo io dico che in realtà \emph{GTA} è un
\textbf{finto MMO}.'' \emph{GTA Online}, soprattutto, è fatto di comunità che fanno
\emph{role play}: c'è chi fa il guidatore delle rapine, perché esiste la meccanica
per farle, e il gruppo lo designa (``tu aspetti fuori, quando usciamo salti su e
portiamo via l'incasso''). Si avvicina a un MMO, e tutto funziona: ``ma funziona
abbastanza? Questa è la mia domanda.'' Basterà a risollevare Rockstar, e
un'industria in una crisi ``abbastanza forte''? ``È un salto nel vuoto. L'unica
compagnia che lo può fare è Rockstar.'' E se non ci riesce, lo si saprà tra anni:
bisognerà calcolare i risultati di \emph{GTA Online} (abbonamenti, vendita di
denaro di gioco, contenuti, macchine), non quelli di \emph{GTA 6}, ``e diffida di chi
ti dice di aver tratto delle conclusioni entro uno o due anni''.

\paragraph{Il nuovo online.} Un altro studente ricorda che con il preordine di
\emph{GTA 6} venivano date \emph{Shark Card} per \emph{GTA 5}. ``Certo: nel momento
in cui annunciano \emph{GTA 6 Online}, tutti smettono di giocare a \emph{GTA 5
Online}.'' La mappa di \emph{GTA 6} sarà la mappa del nuovo online; c'è anche la
possibilità, ipotizza il docente, che lo aggiungano a quello esistente (``prendi
l'aereo da Los Santos e atterri a Vice City'') e ``allora sì che diventa un MMO''.
Di sicuro lo annunceranno solo quando saranno pronti, dopo che la gente avrà
comprato il gioco, per non bruciare i ricavi dell'online attuale. ``Normalmente ti
direi che secondo me non lo fanno; ma Rockstar ha una quantità di soldi così
enorme che potrebbe riuscirci. Vedremo.''

\paragraph{I valori di produzione.} C'è perfino un internet interno, con i social
media del gioco. ``Se non ci giochi, non capisci la quantità di \textbf{valori di
produzione}'': sono tutte cose non difficili in sé, ``ma la cosa principale che
richiedono è soldi, una quantità di soldi enorme''. Chi lo comprerà? Il docente
aspetterà: pare che su console non giri benissimo, il che è comprensibile, perché
la PlayStation~5 è uscita sei anni fa (``stai girando su hardware di sei anni
fa''). La pausa, spiega, serve anche a lui, ``per non perdere la voce parlando sei,
sette ore di seguito''.

\subsection{Piazzare oggetti via codice: le spline e i geometry node}

Uno studente fa una domanda più tecnica: il posizionamento degli oggetti in una
scena \textbf{tramite codice} si fa solo per i giochi procedurali, o anche per le
rifiniture, per essere precisi al pixel?

\paragraph{Le staccionate.} ``Ci sono procedure diverse.'' In Unreal, uno degli usi
delle \textbf{spline} (introdotte la volta precedente,
sezione~\ref{gd11:sec:spline}) sono le \textbf{staccionate}. Si ha l'asset di un
pezzo di staccionata e si deve circondare un quartiere americano, dove ogni casa
ha la sua. Mettere a mano ogni pezzo, a ogni incrocio, lungo ogni curva, richiede
un tempo enorme. In Unreal si può fare un \textbf{Blueprint} che, nel
\textbf{Construction Script}, dice: ``ho questa spline, metti i pezzi di staccionata
lungo la spline''. Trascinando la spline, la staccionata si costruisce da sola.
\begin{itemize}
  \item In \textbf{Blender} queste cose si fanno con i \textbf{geometry node}.
  \item Con \textbf{Unity} spesso si costruisce l'ambiente direttamente in Blender,
        non dentro l'engine (sezione~\ref{gd11:sec:brush}).
\end{itemize}

\paragraph{In editor, non al buio.} La costruzione avviene nel Construction Script
``perché così puoi vedere il risultato in editor; altrimenti stai lavorando al
buio''. Lo studente chiede se, quando si preme Play, venga tutto ricalcolato: si
fa anche quello, ``a seconda delle necessità del progetto'', soprattutto se
l'ambiente si evolve o si trasforma durante il gioco, o se il gioco è
\textbf{procedurale}: \emph{Minecraft} ``non ha un editor dei livelli, ha un
generatore algoritmico''.

% ---------------------------------------------------------------------
\section{Troppe luci: quando Unreal spegne le ombre}
\label{gd14:sec:troppe-luci}

``Riprendiamo. L'ultima ora, giornata pesante.'' Prima un controllo che ci siano
tutti, anche Federica da casa. Poi un'avvertenza legata al discorso della mobilità
e di Lumen: \textbf{Unreal ha la tendenza a non farti sapere quando ti stai tirando
la zappa sui piedi da solo}. Ci sono situazioni in cui Unreal \textbf{disattiva cose
senza dirlo}, e una riguarda proprio le luci.

\paragraph{Le luci che si spengono.} Se si piazzano \textbf{troppe luci capaci di
proiettare ombre, una vicino all'altra}, Unreal comincia a \textbf{spegnerle}. Uno
studente l'ha appena sperimentato: ha messo sei point light. ``Esattamente: si
disattivano perché Unreal\footnote{Nella trascrizione, per un errore di
riconoscimento, ``Harry Potter''.} automaticamente dice: questa roba è una quantità
allucinante di calcoli, per sicurezza la disattivo.'' Sulle luci coinvolte compare
un simbolino \textbf{X} (rosso): le luci magari illuminano ancora l'ambiente, ma
\textbf{smettono di proiettare ombre}.

\paragraph{Non toccare le impostazioni: imparare a risolvere.} Il consiglio, finché
si è studenti, è di \textbf{non andare a cambiare le impostazioni} per aggirare il
limite, perché si rischia di creare grossi problemi ai propri progetti: meglio
imparare come risolverlo. Le soluzioni sono tre:
\begin{enumerate}
  \item cambiare la \textbf{mobilità} a statica (ma, con Lumen attivo, non è detto
        che cambi qualcosa, sezione~\ref{gd14:sec:mobilita});
  \item togliere il \textbf{Cast Shadows};
  \item \textbf{distanziare} le luci, in modo che non si sovrappongano.
\end{enumerate}

\subsection{Il Cast Shadows}

Il docente lo mostra in \emph{Medieval Village}. Mette una point light accanto
all'insegna di un negozio: la luce proietta l'ombra dell'insegna sul muro. Nei
\emph{Details}, invece di scorrere fino in fondo, si usa la casella di
\textbf{ricerca} (``che è il modo più rapido, se sapete già come si chiamano le
cose''): si scrive \emph{shadow} e compare \textbf{Cast Shadows}. Togliendo la
spunta, resta l'illuminazione, ma \textbf{la luce non proietta più ombre}.

\paragraph{Il negozio importante.} Togliere il Cast Shadows ``torna utile per molte
cose''. Supponiamo che in questo livello ci sia un \textbf{negozio importante}, quello
dove il giocatore compra l'arma segreta\footnote{Il nome dell'arma inventata dal
docente per l'esempio è incomprensibile nella trascrizione.}, e che la zona sia
troppo buia: il negozio dovrebbe essere più notevole. L'istinto è prendere una point
light e piazzarla lì. ``E questo è il tipico aspetto che hanno molti \textbf{giochi
studenteschi}: ombre sparate ovunque.'' La luce, messa per far notare che lì c'è
qualcosa di importante, proietta ombre di ogni minima imperfezione, anche solo
delle piante, ombre molto brutte sulla parete. Si fa allora quello che si è visto in
\emph{Lies of P} con la lampada (sezione~\ref{gd14:sec:liesofp}): \textbf{via il Cast
Shadows}. Un lato dell'insegna è così illuminato che si legge, senza il fastidio
fortissimo delle ombre. Poi si regolano i valori necessari, magari riducendo un po'
l'intensità: nella dimostrazione la luce è ``molto più forte di quanto la farei
normalmente, ma è per renderla visibile''.

\paragraph{La porta da notare.} Altro esempio: il giocatore arriva camminando da una
certa direzione, e si vuole che noti una \textbf{porta} lì davanti. Si mette una point
light, si tolgono le ombre: le ombre non distraggono, e la porta diventa più
visibile da quel punto di vista.
\begin{quote}
Vi ricordate il discorso verosimile e non realistico? Togliere le ombre è
importante, perché se non le togli rovini la verosimiglianza. È un
\textbf{vantaggio che hanno i videogiochi rispetto ai film}: possiamo illuminare
quello che ci serve che sia notato, senza sbatterlo in faccia al giocatore.
\end{quote}
E c'è un vantaggio pratico: \textbf{le luci che non proiettano ombre se ne possono
mettere molte}, a differenza delle altre.

\paragraph{Ambienti, non livelli.} ``E voglio che sia chiara anche questa cosa:
\textbf{noi per ora non creiamo livelli, noi per ora creiamo ambienti}.'' Se in un
ambiente c'è una parte troppo buia, si usa una point light (o una spot light) senza
ombre per alzare un po' il livello di luce in quel punto. Se l'effetto non piace, si
può duplicare la luce e aggiustarla. ``Poi lì va a gusto estetico'': la
dimostrazione è ``molto brutale''; normalmente la luce si farebbe \textbf{più larga,
più soffusa, meno potente}. Ma già così, al giocatore che cammina da quella parte,
``è impossibile che non noti che lì c'è un negozio in cui può entrare. Stiamo dando
\textbf{informazioni pratiche}.''

\paragraph{Quando le ombre servono: distanziare.} Se invece le luci non devono
assolutamente smettere di proiettare ombre, perché sono visivamente legate a
proiettori o lampade e sarebbe strano che il giocatore, passando davanti a una,
facesse ombra e davanti all'altra no, \emph{Lies of P} insegna che basta
\textbf{posizionarle in modo che non si sovrappongano}. Distanziate, il problema è
risolto. ``Ancora: sono cose che i giocatori, se non si fermano a guardarle, non
notano neanche.''

% ---------------------------------------------------------------------
\section{La spot light}
\label{gd14:sec:spot}

La \textbf{spot light}, vista la volta precedente, è ``una point light che fa effetto
torcia'': si dirama sempre da un punto, ma i raggi sono \textbf{limitati all'interno
di un angolo}, un cono (figura~\ref{gd14:fig:tipi}). Simula una torcia, un faro, un
faretto: ``se avete mai usato una torcia in un videogioco horror, è una spot
light''. Come per la point light, si può \textbf{allargare l'origine} (il
\emph{source radius}) perché la luce non nasca da un punto ma da una superficie.
Anche qui ci sono molti valori da toccare, come il \textbf{cono interno} e il
\textbf{cono esterno}: ``vi lascio esplorare da soli; se avete dubbi, chiedetemi
pure''. Non si vede ogni singolo parametro, ``se no non riusciamo a fare tutta la
lezione''.

% ---------------------------------------------------------------------
\section{Medieval Village è giocabile}
\label{gd14:sec:play}

Una piccola curiosità prima di passare oltre: \emph{Medieval Village} \textbf{è
giocabile}. Premendo Play (scorciatoia \textbf{Alt+P}), il computer si blocca per un
attimo, ``perché è una scena professionale di alto livello''. C'è una breve
sequenza iniziale, che il docente salta, e poi si può esplorare il villaggio
``proprio come un videogioco''. Un tasto (\textbf{Q}, nella dimostrazione) toglie
l'interfaccia. È un piccolo ambiente esplorabile, con \textbf{pareti invisibili}
all'esterno, senza nessun tipo di gameplay: ma chi è curioso di vedere come sarebbe
giocarci dentro può farlo.

% ---------------------------------------------------------------------
\section{Leggere una mappa: Interiors Volume 1}
\label{gd14:sec:scanlabs}

\subsection{Aprire la mappa}

Ora tocca al pacchetto di Scan Labs. Si importa nel proprio progetto ``nel modo che
avevo detto'': si prende la cartella dei file già importati, la si mette nella
cartella \texttt{Content} del progetto (dall'interno di Unreal o dal launcher) e si
apre Unreal, o lo si chiude e riapre se era già aperto. ``Ricordatevi di
spacchettarli, di decomprimerli.'' Dentro il pacchetto c'è una mappa chiamata
\textbf{Interiors Volume 1}: il docente dà qualche minuto perché tutti la aprano.

\paragraph{Una mappa vecchia.} È una mappa che il docente usa da tempo come esempio
per gli studenti, fatta con una versione precedente di Unreal, ``quindi mi sa che la
luce non è super adattata, però funziona''. Lumen ``si comporta in una maniera
strana'': ci sono \textbf{macchie di luce che si spostano} sul muro quando si gira
la telecamera. È Lumen che approssima l'illuminazione reale in un ambiente che non
è stato studiato per questa versione di Unreal. ``Tuttavia è ancora perfettamente
funzionale per i nostri scopi, ed è ancora molto bella da vedere.''

\paragraph{Perché queste due mappe.} Perché lasciare agli studenti questa mappa e
\emph{Medieval Village}? Sono gli esempi a cui il docente vuole che si dia
un'occhiata ``per iniziare a \textbf{tarare} i vostri gusti personali e la direzione
di quello che state facendo''. Il percorso fin qui: si è presa la mano con le forme
di base, si sono costruite le prime cose, si è iniziato con lo Starter Content.
\textbf{Adesso si comincia a lavorare con i prop}: si hanno a disposizione cose che
permettono di costruire ambienti ``in una versione semifinale''.
\begin{itemize}
  \item ``Voglio che iniziate a studiare. Voglio che iniziate a
        \textbf{osservare i videogiochi a cui giocate}'': come sono organizzati
        gli ambienti, che cosa fanno, gli spazi, le dimensioni degli spazi,
        ``tutto quello che vi colpisce''.
  \item Questo livello di qualità è l'\textbf{obiettivo ultimo}: ``non mi aspetto
        che lo facciate entro la prossima lezione, ovviamente: è un percorso. Però è
        una buona idea, quando si corre, tenere a mente dov'è il traguardo.''
\end{itemize}

\subsection{Lezione numero uno: occhio non vede, cuore non duole}

``Diamo un'occhiata a questa mappa. Cosa sta succedendo? Chi ha voglia di dare delle
descrizioni? Qualunque cosa. Alleniamoci a guardare.'' Sembra il salone di una
villa, ci sono libri, ``e fino a lì non ci piove''. Ma c'è un'uscita: che cosa
succede, fuori?

\paragraph{Fuori non c'è niente.} Fuori c'è un pezzo di verde e un'immagine sullo
sfondo, ``perché tanto è tutto quello che si vede dall'interno della casa, e va
benissimo''. ``Lezione numero uno'': questa è una mappa pensata per essere vista
\textbf{da un certo punto}, e non ci si sposta più di così. ``\textbf{Occhio non vede,
cuore non duole.}'' Se non serve fare qualcosa, la si lascia fuori dall'inquadratura.
Anche in un gioco ambientato all'esterno, uscendo dai confini con la telecamera
libera dell'editor, si scopre che è così: aggiungere altro sarebbe ``calcoli
sprecati, tempo sprecato''.

\paragraph{Lavorare dalla telecamera.} Il metodo (che si vedrà più avanti): si piazza
una \textbf{telecamera}, la si tiene fissa in un riquadro della viewport, e si
lavora in base a quello che si vede dalla telecamera, riempiendo gli spazi
inquadrati; poi magari si sposta un po' la telecamera per controllare che vada
bene. ``Quindi, numero uno: \textbf{mai sprecare energie}. Ragionate su quello che si
vede e quello che non si vede: se non si vede, non lo fate.'' Con una precisazione:
nei videogiochi la telecamera si muove, ma per ora si fanno solo screenshot.
``Pensiamo ad ambienti, non a livelli.''

\subsection{Numero due: la luce, e il chiaroscuro}

``Cosa si può dire della luce in questa stanza?'' C'è un riflesso; la luce viene da
fuori. La stanza è illuminata prevalentemente da una luce \textbf{direzionale} che
entra attraverso la finestra, secondo regole di \textbf{chiaroscuro} di cui si
parlerà quando si affronterà il lato più artistico delle luci. La luce
``\textbf{dipinge} questi sprazzi che attraversano la stanza'', rendendo più
interessante la \textbf{composizione visiva}.

\paragraph{Caravaggio e il noir.} La luce che entra dalle finestre, e le ombre in
generale, possono essere usate come \textbf{elementi compositivi artistici}.
``Di Caravaggio ne parleremo durante una lezione''. Perché
``questo non è un corso di tutorial di Unreal: questo è un corso di level design'',
e quindi ci saranno momenti in cui si parlerà di Caravaggio, del
\textbf{noir}, della tecnica del \textbf{chiaroscuro}, di teoria artistica di inizio
Novecento: ``è una delle mie lezioni preferite''.

\paragraph{Reflection capture.} Ripremendo G compaiono le icone di tutto ciò che c'è.
Molte sono \textbf{reflection capture} (in altri engine si chiamano
\emph{reflection probe}): ``non preoccupatevene, servono per i riflessi, e i riflessi
sono un discorso a parte''. Precalcolano i riflessi degli oggetti intorno, e sono
legate alle cose \textbf{statiche}. Questo ambiente, infatti, è stato creato pensando
alle \textbf{luci statiche}, mentre ora Unreal lavora con Lumen: la rifinitura delle
luci era pensata per un approccio diverso. È il motivo per cui la stanza sembra un
po' più buia del normale e ha tutto quello sfarfallio.

\subsection{Il pizzico di caos, e i libri}

``Quello che mi interessa è il lato \textbf{compositivo}.''
\begin{itemize}
  \item La \textbf{caoticità} delle scarpe appoggiate sotto il divano, un po'
        storte.
  \item I \textbf{libri}, tutti di dimensioni leggermente diverse tra loro. E
        cliccandoci sopra si scopre che \textbf{non sono individuali}: è
        un'\textbf{unica mesh}, un gruppo di libri in un solo modello. Altrove,
        però, il pacchetto ha anche le mesh dei \textbf{libri singoli}: per chi
        vuole usarli individualmente e, per ottimizzare, hanno fatto anche le
        mesh dei libri a gruppi.
  \item I bicchieri, e così via.
\end{itemize}
``Non ho bisogno di farvi troppi discorsi'': come creazione di interni, la classe
parte ``da un livello abbastanza buono''.

\subsection{Le tre regole dell'osservazione}

Il consiglio più importante della lezione riguarda \textbf{la capacità di
osservare}, che ha tre aspetti (figura~\ref{gd14:fig:osservare}).

\paragraph{Uno: i prodotti finali.} Per lavorare nel mondo dei videogiochi bisogna
saper \textbf{analizzare i videogiochi}, i prodotti finiti. Chi lavora nel cinema
guarda film e capisce come funziona il mondo dei film attorno a sé; chi lavora
nella musica ascolta album per capire come sono fatte le produzioni.

\paragraph{Due: dietro le quinte.} La seconda è poter \textbf{vedere dietro le
quinte}: per questo il docente ha dato questi progetti. Per ora si resta
nell'ambiente del \textbf{realistico moderno} (\emph{Medieval Village} è stato dato
``un po' per sfizio''), così che tutti abbiano lo stesso punto di partenza e più o
meno la stessa idea di cosa si sta facendo.

\paragraph{Tre: le persone al tuo livello.} La terza è ``quella più importante''. Ci
sono necessariamente grosse differenze tra i task di una persona e quelli di
un'altra: ``non esiste una classe in cui tutti fanno i task nella stessa maniera''.
È normale: il level design è una \textbf{specializzazione}, e il docente non pretende
che tutti siano allo stesso livello, perché chi non diventerà level designer deve
sapere come lavora un level designer, ma saperlo fare perfettamente ``potrebbe non
essere la cosa che più gli torna utile''. Però vuole che ciascuno \textbf{guardi i
risultati degli altri} (è uno dei motivi per cui il primo anno i lavori si
correggono tutti insieme) e si chieda se è nella media o no. Scan Labs e
\emph{Medieval Village} sono fatti da gente espertissima: ``è chiaro che voi non
siete a livello di fare quelle robe. Ma \textbf{qual è il livello a cui potreste
farle}? Lo scoprite guardandovi attorno'', guardando i compagni di classe e gli
altri studenti.

\begin{quote}
Regola 1: osservare \textbf{i prodotti finali} della tua industria. Regola 2:
osservare \textbf{che cosa succede dietro le quinte}. Regola 3: osservare
\textbf{che cosa è in grado di fare la gente al tuo livello}. La terza è quella che ti
radica nella realtà in cui ti trovi.
\end{quote}

\begin{figure}[htbp]
\centering
\begin{tikzpicture}[
  gd14rule/.style={draw, rounded corners=3pt, thick, text width=3.6cm,
                   minimum height=2.1cm, align=center, font=\scriptsize},
]
  \node[gd14rule, fill=blue!8] (r1) at (0,0)
    {\textbf{1. I prodotti finali}\\[2pt] giocare e analizzare i giochi finiti,
     come chi fa cinema guarda film};
  \node[gd14rule, fill=green!8] (r2) at (4.4,0)
    {\textbf{2. Dietro le quinte}\\[2pt] aprire progetti professionali
     (\emph{Medieval Village}, Scan Labs) e smontarli};
  \node[gd14rule, fill=orange!12] (r3) at (8.8,0)
    {\textbf{3. Le persone al tuo livello}\\[2pt] guardare i task dei compagni:
     che cosa \emph{posso} fare, adesso?};
  \node[font=\scriptsize\itshape, align=center] at (4.4,-1.65)
    {dove sta il traguardo (1--2) e dove sei tu (3): la terza ti radica nella
     realtà};
  \draw[-{Stealth}, thick] (r1) -- (r2);
  \draw[-{Stealth}, thick] (r2) -- (r3);
\end{tikzpicture}
\caption{Le tre regole dell'osservazione proposte dal docente. Ricostruzione
schematica.}
\label{gd14:fig:osservare}
\end{figure}

\subsection{Rafforzare le luci}

``Qualcuno ha altre osservazioni?'' Uno studente nota dei punti in cui la luce sembra
più forte del previsto e chiede se lì ci siano delle luci aggiunte. ``Normalmente
sì.'' Si farebbe proprio questo: il docente lo chiama \textbf{rafforzare le luci}
(``un'altra delle cose che hanno un sacco di nomi diversi''). Si mette una
\textbf{leggerissima} point light o spot light nei punti in cui sembra esserci troppa
ombra, e si toglie il Cast Shadows, così non si nota
(sezione~\ref{gd14:sec:troppe-luci}). ``È una buona osservazione.'' In questo caso,
però, quelle icone sono reflection capture, e quello che si vede è un'altra cosa: la
\textbf{luce riflessa}, simulata da Lumen.

% ---------------------------------------------------------------------
\section{Lumen, la luce riflessa e i suoi artefatti}
\label{gd14:sec:lumen}

\paragraph{La luce rimbalzata.} La luce colpisce gli oggetti nella stanza, e gli
oggetti emettono a loro volta una \textbf{luce riflessa}, anche se non sono
riflettenti: è luce ambientale. Con Lumen, se si punta una luce su un \textbf{oggetto
rosso}, attorno all'oggetto si crea un \textbf{alone di luce leggermente rossa},
``come succede nella realtà, perché tutti gli oggetti, escluso il nero puro,
emettono luce riflessa quando vengono colpiti dalla luce''. (La dimostrazione in
classe riesce a metà: la luce usata è troppo debole perché l'effetto si veda bene.)

\paragraph{Gli artefatti.} Ed ecco ``una piccola recriminazione personale'': in
certi giochi si nota molto la differenza tra chi sa usare l'illuminazione di Unreal
e chi no. Quelli che si vedono nella mappa sono \textbf{artefatti} della luce: cose
calcolate che svaniscono e appaiono; il \textbf{ghosting}, l'immagine fantasma che
resta quando si sposta un oggetto; un \textbf{rumore} di fondo fatto di puntini e
pallini. Sono l'effetto di un'illuminazione dinamica sempre accesa che \textbf{non è
ray tracing}: se Lumen fosse in tempo reale, sarebbe ray tracing, e il ray tracing
richiede hardware dedicato perché è complicato da calcolare. Lumen invece lo
calcola \textbf{spalmato su più frame}, ``per semplificarsi la vita'', ed è ``super
rumoroso'' perché così deve fare meno calcoli. ``Se avete giocato al remake del
nuovo \emph{Halo}, l'avete notato quasi di sicuro, magari con le opzioni più
basse.''\footnote{Così nella trascrizione; il riferimento esatto non è
ricostruibile con certezza.}

\paragraph{Un compromesso.} Il vantaggio è che si possono spostare oggetti, avere
ombre e così via \textbf{senza doverci pensare}: niente \emph{baking}, e niente
\emph{baking} da rifare ogni volta che si modifica qualcosa, come succede con
l'illuminazione statica; e queste approssimazioni sono relativamente leggere.
L'illuminazione di oggi è basata sull'illuminazione dinamica; ci sono giochi che
\textbf{richiedono obbligatoriamente il ray tracing}, come l'ultimo \emph{Doom}.
``Io personalmente lo trovo bruttissimo da vedere, questo effetto.'' Ma è un tratto
comune dei giochi fatti con Unreal, ``una di quelle cose che contribuiscono al fatto
che Unreal ha l'aspetto di Unreal'', e che si nota molto se gli studi non lavorano
per compensarlo.

\paragraph{Chi lo fa benissimo.} Un esempio di gioco che lo fa benissimo è
\emph{Silent Hill}\footnote{Nella trascrizione ``Silent Till Have''; quasi certamente
uno dei recenti \emph{Silent Hill} realizzati con Unreal Engine 5 (il remake di
\emph{Silent Hill 2} o \emph{Silent Hill f}): non è chiaro quale.}: ``non sembra fatto
con Unreal, da quanta attenzione ci hanno dedicato alla luce''. Essendo un horror,
alla luce hanno dedicato tantissima attenzione.

\paragraph{Una parentesi con uno studente.} Ne nasce una breve conversazione. Su un
gioco Nintendo per Switch~2, il docente commenta che ``i laboratori di ricerca e
sviluppo di Nintendo fanno delle robe incredibili''. Si parla di aziende che fanno
giochi per mettere alla prova la propria tecnologia: a Nvidia ``non interessa fare
videogiochi'', ma la \textbf{Blender Foundation} li fa, per testare sul campo il
proprio programma di modellazione e scovarne i problemi; \emph{Dogwalk} è il loro
ultimo gioco, ``molto carino'', e sono gratis.

\paragraph{Lumen non è ray tracing.} ``Attenzione: Lumen è un'approssimazione molto
approssimata del ray tracing.'' Il paragone serve a dare l'idea, ma Lumen si chiama
\textbf{illuminazione dinamica} (illuminazione globale dinamica). Il ray tracing si
può fare via \textbf{hardware} o via \textbf{software}: se la scheda grafica ha parti
dedicate al ray tracing lo si fa molto più facilmente; via software si può, ``ma ti
gira malissimo''. Dire che Lumen è ray tracing software ``non è esattamente
corretto'', ma lo scopo è lo stesso: \textbf{simulare un'illuminazione dinamica che
considera tutte le fonti di luce}. Il ray tracing lo fa in modo ``più o meno
perfetto'', anche se in realtà anche lì c'è molto upscaling e molta
approssimazione, perché il ray tracing reale è impossibile: calcolare per ogni pixel
non solo tutte le fonti di luce, ma come la luce rimbalza creando fonti di luce a
cascata, sarebbe una quantità di calcoli irraggiungibile.

\paragraph{Le prime schede.} ``Io mi ricordo quando sono uscite le prime schede'':
accendere il ray tracing faceva ``implorare pietà'' al computer, il gioco girava a
quindici fps, ed era un ray tracing ``molto banale''. Il docente era presente, per
lavoro, all'evento in cui Nvidia presentò per la prima volta le schede con ray
tracing, a Colonia.\footnote{Nella trascrizione ``alla GDC a Colonia'': la fiera di
Colonia è la Gamescom; la presentazione delle prime schede RTX avvenne lì nel 2018.}

\paragraph{Per ora, Lumen acceso.} ``Non divaghiamo troppo.'' Per ora si tiene Lumen
acceso e ci si fa una ragione dei suoi difetti: niente modifiche alla mobilità delle
luci. F11 per tornare alla finestra normale.

\paragraph{Approfondimento: ombra e penombra secondo Totten.} Totten distingue due
condizioni di luce utili al level designer. La \textbf{penombra} (\emph{shade}) è la
luce diffusa che entra attraverso tende, schermi o vetrate: era una caratteristica
dell'architettura gotica (la \emph{lux nova}), ed è un caposaldo dell'horror; poiché
un'informazione visiva incompleta invoglia a esplorare, la penombra crea attrattiva
e ambiguità (lo spazio è sicuro o pericoloso?). L'\textbf{ombra} (\emph{shadow}) è
invece la condizione creata quando un oggetto copre la sorgente: può nascondere chi
vi si trova (gli sviluppatori di \emph{Splinter Cell} parlavano di
\emph{shadowspace}) e può creare un \textbf{ritmo}, come le ombre delle finestre
proiettate perpendicolarmente al percorso del giocatore, che lo ``tirano'' avanti. È
lo stesso uso compositivo degli sprazzi di luce della mappa di Scan Labs: la luce che
entra dalla finestra non illumina soltanto, \emph{disegna} lo spazio.

% ---------------------------------------------------------------------
\section{I materiali sulle luci}
\label{gd14:sec:lightfn}

Un'ultima nota, ``più che altro come curiosità'', per chi ha già fatto corsi extra
ed è ``molto lanciato'': si possono \textbf{aggiungere materiali alle luci}. Che cosa
vuol dire? Che \textbf{una luce proietta un'immagine}.
\begin{itemize}
  \item Il docente l'aveva usato in un prototipo in cui i \textbf{tutorial} erano
        proiettati sulle pareti da degli oggetti, che mostravano che cosa bisognava
        fare.
  \item Se in un videogioco la \textbf{torcia} non è luce pura ma ha quella tipica
        forma circolare delle torce, è perché la proiettano così.
  \item In \emph{Control}, quando un proiettore manda un video su una parete e gli
        si spara, il proiettore rotola e il video, proiettato in tempo reale, si
        sposta lungo la stanza su tutto ciò che c'è intorno.
\end{itemize}
È la stessa logica: \textbf{un materiale, anche animato, assegnato alla luce}, che
permette di proiettare un'immagine su una parete; e se ci passa davanti il
giocatore, lascia l'ombra e tutto quanto. Sulla slide ci sono le istruzioni per
impostare il materiale, ``ma cosa sono i materiali lo vedremo in futuro'': ``non mi
interessa sinceramente tanto che lo facciate''.

% ---------------------------------------------------------------------
\section{Il task: una stanza illuminata con uno scopo}
\label{gd14:sec:task}

Il docente riscrive sul momento la consegna, ``perché in realtà è sbagliata''.

\begin{quote}
Create \textbf{un'altra stanza}, sullo stile di quelle fatte finora, con le stesse
regole usate finora. \textbf{Non usate} il progetto medievale. Potete usare
\textbf{Scan Labs} e lo \textbf{Starter Content}, ma \textbf{non riutilizzate la
stanza di Scan Labs}: costruite l'ambiente da zero, con la geometria per i muri e i
materiali dello Starter Content. Questa volta \textbf{ponete attenzione alle
luci}: usatele in maniera \textbf{logica}, secondo uno \textbf{scopo preciso}, e
\textbf{spiegate lo scopo} con cui le avete usate. \emph{Make it a significant
room}.
\end{quote}

\paragraph{Le pareti di Scan Labs.} ``Questa è una cosa che mi dicono spesso'': Scan
Labs ha una stanza le cui pareti sono già fatte a forma di stanza. Il docente non
vuole che la si riutilizzi, ``a meno che non vogliate proprio cambiarla in maniera
generale''; preferisce che si crei un ambiente da zero.

\paragraph{Volutamente vago.} ``È ovvio che lo sto dicendo in una maniera molto vaga'':
il docente vuole lasciare la libertà di interpretare l'obiettivo come si preferisce.
Per luci si intendono point light, spot light, ``qualunque cosa vi venga in mente.
Se avete un'idea, tiratela giù.''

\paragraph{Una stanza significativa.} La regola già detta a inizio lezione
(sezione~\ref{gd14:sec:revisione}) diventa esplicita nella consegna: ``vi ho lasciato
la libertà per il task precedente; \textbf{d'ora in poi ci concentriamo su risultati
pseudo-realistici}, invece di fare quello che ci pare, per allenarci ad avere un
obiettivo preciso. Poi a un certo punto le regole cambieranno.''

\paragraph{Cominciare in classe.} Manca mezz'ora alla fine: il docente fa cominciare
subito il lavoro, così che eventuali problemi di importazione e domande si risolvano
in aula, ``e non devo rispondere a tre persone che mi fanno la stessa domanda su
Discord''. Dopo dieci minuti fa il giro dei banchi. (Il giro, compresa la riparazione
dei materiali di uno studente, recuperando a mano la texture \emph{diffuse} giusta,
è in gran parte incomprensibile nella registrazione.)

% ---------------------------------------------------------------------
\section{Le domande sul task}
\label{gd14:sec:domande}

\subsection{Fab, no; le reference, sì}

\paragraph{Fab.} Prima domanda: valgono le regole della volta scorsa, cioè niente
\textbf{Fab}, solo i pacchetti dati? ``Sì, sempre questa regola: per ora non stiamo
usando Fab.'' A un certo punto il docente dirà di scegliere che cosa scaricare da
Fab, ma per ora ci si concentra su quello che è stato dato. Quello che vuole evitare
è: ``vai su Fab, trovi il pacchetto della stanza, rifai la stanza di Fab''.

\paragraph{La stanza di Lynch.} Una studentessa vorrebbe ricreare come punto di
partenza una stanza di un film o di una serie TV (``una stanza di Lynch, che ne
so''), con lo stesso pavimento: se trova l'immagine su internet e ne ricava lei il
materiale, vale come averlo scaricato? ``No. Per queste task potete usare solo i
pacchetti che vi ho dato: lo Starter Content e Scan Labs. Se volete imitare qualcosa,
se volete usare una reference, siete assolutamente liberi.''

\paragraph{Arrangiarsi con quello che si ha.} ``Lo scopo di queste lezioni è
imparare ad arrangiarsi con quello che si ha.'' Ed è una cosa che il docente aveva già
detto individualmente a qualcuno, ma vuole ripeterla a tutti:
\begin{quote}
Non voglio che vi abituiate a lavorare dicendo: mi serve questa cosa, vado a prenderla
su Fab e la butto dentro. Una delle capacità più importanti che dovete sviluppare
(e queste lezioni sono un momento importante per farlo, in uno spazio sicuro,
tranquillo e pensato per allenarsi) è \textbf{la capacità di arrangiarvi con quello
che avete}.
\end{quote}
Si devono creare, o ricreare se si lavora da reference, spazi in base a ciò che si
ha a disposizione. ``Altrimenti la vostra abilità principale diventa cercare su Fab,
e nel momento in cui andate a lavorare su Unity siete fregati, detto proprio terra
terra.'' Il docente non vuole che ci si abitui a lavorare solo con Unreal, ``perché
non è detto che lavorerete solo con Unreal'' e, a livello personale, perché non gli
piace che ci sia ``un monopolio così fisso'' sull'engine standard dei tripla A.

\subsection{Si parte dalla luce o dall'ambiente?}

Un'altra domanda: di solito si costruisce prima la scena e poi si piazzano le luci
dove rendono meglio, oppure si parte già dall'idea di luce e si costruisce l'ambiente
attorno? ``È nato prima l'uovo o la gallina? La risposta è: che idea hai?'' Se l'idea
si basa su un gioco di luce, si parte dalla luce; se si basa su un ambiente (``come
lo illumino?''), si parte dall'ambiente. In queste condizioni non c'è una vera
risposta: quando si parlerà di layout, concettualizzazione e blockout ci saranno
procedure precise, ma per ora si fanno ambienti per allenarsi. ``Se visualizzi la
luce che penetra dalle sbarre di una finestra, parti dalla finestra con le sbarre e
poi costruisci tutto; se ti piace la scrivania, parti dalla scrivania. Basta che
trovi un \textbf{punto di partenza} a cui ancorarti.''

\subsection{Modificare un materiale dello Starter Content}

``Ultimissima cosa'': si può modificare un po' un materiale che è già nello Starter
Content? La lezione sui materiali non c'è ancora stata, ``ma se è un materiale che
già c'è nelle cose che vi ho dato, allora puoi andare a toccarlo, dai. Mi sembra
esagerato dire che non puoi neanche toccarlo.'' Chi non ne ha voglia, non lo faccia:
ci sarà una lezione dedicata. ``Quello te lo concedo.''

\subsection{Le luci al neon sugli scaffali}

Durante il giro, uno studente discute con il docente un'idea: delle \textbf{luci al
neon} lungo gli scaffali, tipiche di certe case di alto livello. Il docente osserva
che, in un ambiente moderno, una cosa del genere va \textbf{contestualizzata}: o i
neon sono colorati e diventano un elemento caratterizzante, oppure diventano una
fascia di luce che sottolinea l'elemento principale della stanza. Per una fascia di
luce lunga potrebbe servire una \textbf{rect light}. Soprattutto, deve essere una
stanza \textbf{funzionale}: ``vissuta'', credibile. Non deve essere un tipo preciso
di stanza (``se mi fai un corridoio, va bene''), basta che sia un
\textbf{interno} contestualizzato.

% ---------------------------------------------------------------------
\section{La lezione spostata e la chiusura}
\label{gd14:sec:chiusura}

\paragraph{La lezione della settimana prossima.} Un avviso: la lezione della
settimana successiva va \textbf{spostata}, perché il docente ha una visita medica il
mercoledì pomeriggio. La data la comunicherà la segreteria (``dovrebbe essere
martedì pomeriggio''), e sarà quasi certamente \textbf{online} (``nove su dieci sarà
da remoto''), perché non c'è un'aula disponibile. In ogni caso tutte le lezioni sono
\textbf{registrate}, così chi ne perde una può recuperarla; le registrazioni vengono
cancellate alla fine dell'anno. ``Vi chiedo scusa del disservizio.''

\paragraph{La chiusura.} ``Ok, allora ci siamo. La lezione si conclude qui. Avete le
vostre task. Se avete bisogno, sapete dove trovarmi; se avete curiosità, sapete dove
trovarmi.'' E, nel caso qualcuno avesse dubbi: ``non mi prendo le ferie per giocare a
\emph{GTA}''. C'è anche uno sciopero annunciato tra giovedì e venerdì: ``vedremo cosa
succede''. ``Ciao Federica. Ci vediamo la settimana prossima.''

% ---------------------------------------------------------------------
\section*{Punti chiave del capitolo}
\addcontentsline{toc}{section}{Punti chiave del capitolo}

\begin{itemize}
  \item \emph{GTA 6} sarà un successo enorme, ma la vera domanda è se basterà: il
        suo scopo è lanciare il nuovo \emph{GTA Online}, e la risposta arriverà
        tra anni. Solo Rockstar può fare ciò che fa Rockstar, anche solo per il
        nome. \emph{Fortnite} non è un gioco ma una piattaforma: uno spazio sociale
        tenuto vivo da contenuti continui, con meccaniche da gioco d'azzardo, che
        vuole diventare ``il mercato'' come Roblox.
  \item \emph{Medieval Village} è un \textbf{progetto} (si apre con doppio clic,
        scegliendo la 5.8); \emph{Scan Labs Interiors} è un \textbf{pacchetto di
        asset} già importati (va decompresso e messo nella cartella
        \texttt{Content}).
  \item Dalla revisione: creare \textbf{spazi coerenti}, non forme a caso; niente
        avanguardia finché non si hanno le basi; \textbf{partire dalla realtà} (una
        reference) e modificarla dà molti più dettagli; il \textbf{manichino} nella
        scena tiene le proporzioni; ciò che conta è che un oggetto sia
        \textbf{leggibile}; un pizzico di caos rende gli ambienti vissuti; con i
        mezzi che mancano ci si arrangia (assi di legno invece di tronchi).
  \item Le viste \textbf{Top/Bottom/Left/Right} sono ortografiche e passano in
        \emph{Unlit}. Sono \textbf{le ombre a dare la profondità}. Il terreno usa
        livelli di dettaglio in base alla distanza. Al primo anno,
        \textbf{dell'ottimizzazione fregatevene} (con qualche asterisco).
  \item Genitore e figlio sono relazioni tra \textbf{actor}; un actor è fatto di
        \textbf{componenti}. Select All Descendants / Immediate Children.
  \item \textbf{PBR}: un metodo di rendering che approssima l'interazione
        realistica della luce con gli oggetti. Un \textbf{materiale} è una scatola
        di texture e numeri; una \textbf{texture} è un'immagine. \textbf{Albedo}
        (= color, base color, diffuse) è solo il colore, ed è ciò che si vede in
        \emph{Unlit}; la \textbf{roughness} sparpaglia i rimbalzi e spegne i
        riflessi (anche con una mappa in scala di grigi: i bordi rovinati); lo
        \textbf{specular} è la riflettenza, e interagisce con la roughness. Gli
        \textbf{specular highlight} si spostano con il punto di vista.
  \item \textbf{Mobilità}: Static (baking in lightmap, leggerissima, niente si
        muove, va buildata con \emph{Build Lighting Only}), Stationary (luce fissa,
        ombre dinamiche), Movable (tutto dinamico, la più costosa). Con
        \textbf{Lumen} tutto è dinamico, approssimato tra un frame e l'altro (da qui
        sfarfallii e rumore). Un valore che ``non fa niente'' può valere solo per
        un tipo di mobilità: leggere il tooltip.
  \item \textbf{Directional light}: raggi paralleli, una sola per livello, conta
        solo la rotazione (la si può spostare dove fa comodo). La notte è il giorno
        tinto di blu. \textbf{Sky Light}: la luce diffusa dall'atmosfera, cioè ciò
        che c'è nelle ombre: cambiandone il colore si cambia l'atmosfera.
  \item \textbf{Verosimiglianza, non realismo}: in \emph{Lies of P} la luce del
        barile è sopra il barile, la lampada non proietta la propria ombra, la luce
        del corridoio sta al centro e la lampada la ``motiva''. Non si mette la luce
        dentro la lampada: si illumina come su un set, per lo scopo della scena.
  \item Troppe luci con ombre vicine: Unreal le \textbf{spegne in silenzio}
        (simbolo X). Soluzioni: statica, \textbf{Cast Shadows} disattivato,
        distanziarle. Le luci senza ombre servono a \textbf{rafforzare} le zone
        buie e a far notare negozi e porte. Per ora si creano \textbf{ambienti, non
        livelli}.
  \item Leggere una mappa: \textbf{ciò che non si vede non si costruisce}; la luce
        dalla finestra e il \textbf{chiaroscuro} come composizione; le reflection
        capture; i libri come mesh unica. Le \textbf{tre regole dell'osservazione}:
        i prodotti finali, il dietro le quinte, il livello delle persone attorno a
        te.
  \item Lumen non è ray tracing: è un'illuminazione globale dinamica approssimata,
        con artefatti (ghosting, rumore) che danno a Unreal ``l'aspetto di
        Unreal''. Le luci possono proiettare materiali (torce, proiettori).
  \item Task: \textbf{una nuova stanza significativa}, pseudo-realistica, con Scan
        Labs e Starter Content (niente Fab, niente progetto medievale, niente stanza
        di Scan Labs riciclata), con \textbf{luci usate per uno scopo} e la
        spiegazione di quello scopo. Si parte da ciò a cui ci si può ancorare: la
        luce o l'ambiente.
\end{itemize}
